% Options for packages loaded elsewhere
% Options for packages loaded elsewhere
\PassOptionsToPackage{unicode}{hyperref}
\PassOptionsToPackage{hyphens}{url}
\PassOptionsToPackage{dvipsnames,svgnames,x11names}{xcolor}
%
\documentclass[
  letterpaper,
  DIV=11,
  numbers=noendperiod]{scrreprt}
\usepackage{xcolor}
\usepackage{amsmath,amssymb}
\setcounter{secnumdepth}{5}
\usepackage{iftex}
\ifPDFTeX
  \usepackage[T1]{fontenc}
  \usepackage[utf8]{inputenc}
  \usepackage{textcomp} % provide euro and other symbols
\else % if luatex or xetex
  \usepackage{unicode-math} % this also loads fontspec
  \defaultfontfeatures{Scale=MatchLowercase}
  \defaultfontfeatures[\rmfamily]{Ligatures=TeX,Scale=1}
\fi
\usepackage{lmodern}
\ifPDFTeX\else
  % xetex/luatex font selection
\fi
% Use upquote if available, for straight quotes in verbatim environments
\IfFileExists{upquote.sty}{\usepackage{upquote}}{}
\IfFileExists{microtype.sty}{% use microtype if available
  \usepackage[]{microtype}
  \UseMicrotypeSet[protrusion]{basicmath} % disable protrusion for tt fonts
}{}
\makeatletter
\@ifundefined{KOMAClassName}{% if non-KOMA class
  \IfFileExists{parskip.sty}{%
    \usepackage{parskip}
  }{% else
    \setlength{\parindent}{0pt}
    \setlength{\parskip}{6pt plus 2pt minus 1pt}}
}{% if KOMA class
  \KOMAoptions{parskip=half}}
\makeatother
% Make \paragraph and \subparagraph free-standing
\makeatletter
\ifx\paragraph\undefined\else
  \let\oldparagraph\paragraph
  \renewcommand{\paragraph}{
    \@ifstar
      \xxxParagraphStar
      \xxxParagraphNoStar
  }
  \newcommand{\xxxParagraphStar}[1]{\oldparagraph*{#1}\mbox{}}
  \newcommand{\xxxParagraphNoStar}[1]{\oldparagraph{#1}\mbox{}}
\fi
\ifx\subparagraph\undefined\else
  \let\oldsubparagraph\subparagraph
  \renewcommand{\subparagraph}{
    \@ifstar
      \xxxSubParagraphStar
      \xxxSubParagraphNoStar
  }
  \newcommand{\xxxSubParagraphStar}[1]{\oldsubparagraph*{#1}\mbox{}}
  \newcommand{\xxxSubParagraphNoStar}[1]{\oldsubparagraph{#1}\mbox{}}
\fi
\makeatother

\usepackage{color}
\usepackage{fancyvrb}
\newcommand{\VerbBar}{|}
\newcommand{\VERB}{\Verb[commandchars=\\\{\}]}
\DefineVerbatimEnvironment{Highlighting}{Verbatim}{commandchars=\\\{\}}
% Add ',fontsize=\small' for more characters per line
\usepackage{framed}
\definecolor{shadecolor}{RGB}{241,243,245}
\newenvironment{Shaded}{\begin{snugshade}}{\end{snugshade}}
\newcommand{\AlertTok}[1]{\textcolor[rgb]{0.68,0.00,0.00}{#1}}
\newcommand{\AnnotationTok}[1]{\textcolor[rgb]{0.37,0.37,0.37}{#1}}
\newcommand{\AttributeTok}[1]{\textcolor[rgb]{0.40,0.45,0.13}{#1}}
\newcommand{\BaseNTok}[1]{\textcolor[rgb]{0.68,0.00,0.00}{#1}}
\newcommand{\BuiltInTok}[1]{\textcolor[rgb]{0.00,0.23,0.31}{#1}}
\newcommand{\CharTok}[1]{\textcolor[rgb]{0.13,0.47,0.30}{#1}}
\newcommand{\CommentTok}[1]{\textcolor[rgb]{0.37,0.37,0.37}{#1}}
\newcommand{\CommentVarTok}[1]{\textcolor[rgb]{0.37,0.37,0.37}{\textit{#1}}}
\newcommand{\ConstantTok}[1]{\textcolor[rgb]{0.56,0.35,0.01}{#1}}
\newcommand{\ControlFlowTok}[1]{\textcolor[rgb]{0.00,0.23,0.31}{\textbf{#1}}}
\newcommand{\DataTypeTok}[1]{\textcolor[rgb]{0.68,0.00,0.00}{#1}}
\newcommand{\DecValTok}[1]{\textcolor[rgb]{0.68,0.00,0.00}{#1}}
\newcommand{\DocumentationTok}[1]{\textcolor[rgb]{0.37,0.37,0.37}{\textit{#1}}}
\newcommand{\ErrorTok}[1]{\textcolor[rgb]{0.68,0.00,0.00}{#1}}
\newcommand{\ExtensionTok}[1]{\textcolor[rgb]{0.00,0.23,0.31}{#1}}
\newcommand{\FloatTok}[1]{\textcolor[rgb]{0.68,0.00,0.00}{#1}}
\newcommand{\FunctionTok}[1]{\textcolor[rgb]{0.28,0.35,0.67}{#1}}
\newcommand{\ImportTok}[1]{\textcolor[rgb]{0.00,0.46,0.62}{#1}}
\newcommand{\InformationTok}[1]{\textcolor[rgb]{0.37,0.37,0.37}{#1}}
\newcommand{\KeywordTok}[1]{\textcolor[rgb]{0.00,0.23,0.31}{\textbf{#1}}}
\newcommand{\NormalTok}[1]{\textcolor[rgb]{0.00,0.23,0.31}{#1}}
\newcommand{\OperatorTok}[1]{\textcolor[rgb]{0.37,0.37,0.37}{#1}}
\newcommand{\OtherTok}[1]{\textcolor[rgb]{0.00,0.23,0.31}{#1}}
\newcommand{\PreprocessorTok}[1]{\textcolor[rgb]{0.68,0.00,0.00}{#1}}
\newcommand{\RegionMarkerTok}[1]{\textcolor[rgb]{0.00,0.23,0.31}{#1}}
\newcommand{\SpecialCharTok}[1]{\textcolor[rgb]{0.37,0.37,0.37}{#1}}
\newcommand{\SpecialStringTok}[1]{\textcolor[rgb]{0.13,0.47,0.30}{#1}}
\newcommand{\StringTok}[1]{\textcolor[rgb]{0.13,0.47,0.30}{#1}}
\newcommand{\VariableTok}[1]{\textcolor[rgb]{0.07,0.07,0.07}{#1}}
\newcommand{\VerbatimStringTok}[1]{\textcolor[rgb]{0.13,0.47,0.30}{#1}}
\newcommand{\WarningTok}[1]{\textcolor[rgb]{0.37,0.37,0.37}{\textit{#1}}}

\usepackage{longtable,booktabs,array}
\newcounter{none} % for unnumbered tables
\usepackage{calc} % for calculating minipage widths
% Correct order of tables after \paragraph or \subparagraph
\usepackage{etoolbox}
\makeatletter
\patchcmd\longtable{\par}{\if@noskipsec\mbox{}\fi\par}{}{}
\makeatother
% Allow footnotes in longtable head/foot
\IfFileExists{footnotehyper.sty}{\usepackage{footnotehyper}}{\usepackage{footnote}}
\makesavenoteenv{longtable}
\usepackage{graphicx}
\makeatletter
\newsavebox\pandoc@box
\newcommand*\pandocbounded[1]{% scales image to fit in text height/width
  \sbox\pandoc@box{#1}%
  \Gscale@div\@tempa{\textheight}{\dimexpr\ht\pandoc@box+\dp\pandoc@box\relax}%
  \Gscale@div\@tempb{\linewidth}{\wd\pandoc@box}%
  \ifdim\@tempb\p@<\@tempa\p@\let\@tempa\@tempb\fi% select the smaller of both
  \ifdim\@tempa\p@<\p@\scalebox{\@tempa}{\usebox\pandoc@box}%
  \else\usebox{\pandoc@box}%
  \fi%
}
% Set default figure placement to htbp
\def\fps@figure{htbp}
\makeatother


% definitions for citeproc citations
\NewDocumentCommand\citeproctext{}{}
\NewDocumentCommand\citeproc{mm}{%
  \begingroup\def\citeproctext{#2}\cite{#1}\endgroup}
\makeatletter
 % allow citations to break across lines
 \let\@cite@ofmt\@firstofone
 % avoid brackets around text for \cite:
 \def\@biblabel#1{}
 \def\@cite#1#2{{#1\if@tempswa , #2\fi}}
\makeatother
\newlength{\cslhangindent}
\setlength{\cslhangindent}{1.5em}
\newlength{\csllabelwidth}
\setlength{\csllabelwidth}{3em}
\newenvironment{CSLReferences}[2] % #1 hanging-indent, #2 entry-spacing
 {\begin{list}{}{%
  \setlength{\itemindent}{0pt}
  \setlength{\leftmargin}{0pt}
  \setlength{\parsep}{0pt}
  % turn on hanging indent if param 1 is 1
  \ifodd #1
   \setlength{\leftmargin}{\cslhangindent}
   \setlength{\itemindent}{-1\cslhangindent}
  \fi
  % set entry spacing
  \setlength{\itemsep}{#2\baselineskip}}}
 {\end{list}}
\usepackage{calc}
\newcommand{\CSLBlock}[1]{\hfill\break\parbox[t]{\linewidth}{\strut\ignorespaces#1\strut}}
\newcommand{\CSLLeftMargin}[1]{\parbox[t]{\csllabelwidth}{\strut#1\strut}}
\newcommand{\CSLRightInline}[1]{\parbox[t]{\linewidth - \csllabelwidth}{\strut#1\strut}}
\newcommand{\CSLIndent}[1]{\hspace{\cslhangindent}#1}



\setlength{\emergencystretch}{3em} % prevent overfull lines

\providecommand{\tightlist}{%
  \setlength{\itemsep}{0pt}\setlength{\parskip}{0pt}}



 


\KOMAoption{captions}{tableheading}
\makeatletter
\@ifpackageloaded{bookmark}{}{\usepackage{bookmark}}
\makeatother
\makeatletter
\@ifpackageloaded{caption}{}{\usepackage{caption}}
\AtBeginDocument{%
\ifdefined\contentsname
  \renewcommand*\contentsname{Table of contents}
\else
  \newcommand\contentsname{Table of contents}
\fi
\ifdefined\listfigurename
  \renewcommand*\listfigurename{List of Figures}
\else
  \newcommand\listfigurename{List of Figures}
\fi
\ifdefined\listtablename
  \renewcommand*\listtablename{List of Tables}
\else
  \newcommand\listtablename{List of Tables}
\fi
\ifdefined\figurename
  \renewcommand*\figurename{Figure}
\else
  \newcommand\figurename{Figure}
\fi
\ifdefined\tablename
  \renewcommand*\tablename{Table}
\else
  \newcommand\tablename{Table}
\fi
}
\@ifpackageloaded{float}{}{\usepackage{float}}
\floatstyle{ruled}
\@ifundefined{c@chapter}{\newfloat{codelisting}{h}{lop}}{\newfloat{codelisting}{h}{lop}[chapter]}
\floatname{codelisting}{Listing}
\newcommand*\listoflistings{\listof{codelisting}{List of Listings}}
\makeatother
\makeatletter
\makeatother
\makeatletter
\@ifpackageloaded{caption}{}{\usepackage{caption}}
\@ifpackageloaded{subcaption}{}{\usepackage{subcaption}}
\makeatother
\usepackage{bookmark}
\IfFileExists{xurl.sty}{\usepackage{xurl}}{} % add URL line breaks if available
\urlstyle{same}
\hypersetup{
  colorlinks=true,
  linkcolor={blue},
  filecolor={Maroon},
  citecolor={Blue},
  urlcolor={Blue},
  pdfcreator={LaTeX via pandoc}}


\author{}
\date{}
\begin{document}

\renewcommand*\contentsname{Table of contents}
{
\setcounter{tocdepth}{2}
\tableofcontents
}

\bookmarksetup{startatroot}

\chapter*{Introduction}\label{introduction}
\addcontentsline{toc}{chapter}{Introduction}

\markboth{Introduction}{Introduction}

\section*{Background}\label{background}
\addcontentsline{toc}{section}{Background}

\markright{Background}

I am currently enrolled in STATS200Q: Philosophical Foundations of
Statistics with Professor Dennis Sun (technically just finished). It is
a small class, with about twelve people, allowing Professor Sun to dive
into the material while maintaining a level of connection with the
students. I have greatly appreciated this element as it has allowed me
to ask many questions along the way!

In my statistical journey, this class served as a follow-up to STATS118,
which I took in the Fall with Professor Gene Kim. Put simply, that class
was very difficult. It involved a lot of mathematical computation and
required students to pick up patterns and manipulate expressions into
known distributions.

As mentioned in the title, STATS200Q feels much more philosophical. It
is still no doubt a math class, but it opts for less computation than
118 and more answering of ``Why?''. Looking back, many of the strategies
I learned in 118 were ``frequentist''. However, even after taking that
final exam, I would not have been able to explain what that term meant.

In STATS200Q, we have learned about the differences between frequentist
and Bayesian statistics and how they approach problems separately. Each
week throughout the quarter, we had a small homework assignment. This
final book is an extension of that. Each assignment is its own chapter,
and I have additionally added a ``Background'' section to each chapter.
Put together, this should make it feel like a tightly-knit narrative.

This introduction chapter is no different. Sadly, this is the end of the
``Background'' section. Below I will explain the distribution and
dataset I focused my studies on. I hope you enjoy!

NOTE: There are chapters in the book with multiple R blocks. Some R
blocks call on variables that were defined in a separate R block earlier
in that chapter. So, please make sure to run the R code in order for
each chapter.

\section*{My Dataset and
Distribution}\label{my-dataset-and-distribution}
\addcontentsline{toc}{section}{My Dataset and Distribution}

\markright{My Dataset and Distribution}

We were each assigned a distribution to focus our analysis on. I chose
the Negative Binomial distribution, and the variant that measures how
many failures before a successful attempt. I focus my data analysis on a
specific COVID hospital dataset, which keeps track of who infected who
in that hospital. In my dataset, the data with the negative binomial
distribution is the number of people each person infected beyond the
first person (This is because all infectors in the dataset infected at
least one person, so we subtract one to have a healthy amount of zero
values).

Source: The original dataset is from this GitHub link:
\url{https://github.com/linwangidd/covid19_transmissionPairs_China/blob/master/transmission_pairs_covid_v2.csv}.

The PMF for the variant of the Negative Binomial I will be analyzing is:

\[
\begin{aligned}
    P(X=x) &= \binom{x + r- 1}{x} p^r (1 - p)^x \\
\end{aligned}
\]

If you want to do any analysis of the data, feel free to code in the
block below.

\texttt{dat} is the dataset.

\texttt{x} is the number of people each person infected beyond the first
person, and follows a negative binomial distribution.

\begin{Shaded}
\begin{Highlighting}[]
\NormalTok{dat \textless{}{-} read.csv("https://raw.githubusercontent.com/Ryder4real/STATS200Q\_Book/main/transmission\_pairs\_covid\_v2.csv")}
\NormalTok{x \textless{}{-} as.vector(table(dat$infector\_id)) {-} 1}
\NormalTok{length(x)}
\end{Highlighting}
\end{Shaded}

As seen in the code block above, there are 807 patients in my dataset.

Traditionally, in this version of the Negative Binomial distribution,
\(p\) is meant to be the probability of success for each trial and \(r\)
is meant to be the fixed number of successes we are waiting for (which
hypothetically would be \(1\) in this case). However, in the real world,
studies have found that negative binomials can effectively model
hospital datasets where the parameters have slightly altered meanings (I
will mention one such study at the end of this chapter).

In my dataset, \(p\) is the probability of a transmission opportunity
not resulting in an infection (perhaps counter-intuitively, the lower
the \(p\), the more infectious). \(r\) is the dispersion parameter,
meaning it is a measure of how heavily clustered the transmissions are.

\(r\) seems to fit best below \(1\), meaning there are many individuals
who do not pass the virus to anyone beyond the first person. This makes
sense as, from the code below, we can see that \(507\) of the
individuals in the dataset have a value of \(0\), meaning they did not
pass the virus to anyone beyond the first person.

\begin{Shaded}
\begin{Highlighting}[]
\NormalTok{sum(x == 0)}
\end{Highlighting}
\end{Shaded}

It is important to note that the PMF I wrote out earlier includes \(r\)
in a choose function. Technically, this is not allowed, as \(r\) is not
necessarily an integer for my dataset. However, in practice, we can just
use the gamma fraction to replace the choose function whenever
computations are needed.

For some quick context, there is a great 2005 paper by Lloyd-Smith et
al. (\citeproc{ref-lloydsmith2005}{2005}) titled ``Superspreading and
the effect of individual variation on disease emergence'', which found
the negative binomial distribution to be a great fit for a 2002-2004
SARS epidemic in Singapore and Beijing, and their \(\hat{r}\) was
\(0.16\).

As a quick note, for this class, we mainly focused on analyzing one
parameter. So, for many of these chapters, I will fix \(r\), and analyze
\(p\).

\bookmarksetup{startatroot}

\chapter{MLE \& Posterior Mean}\label{mle-posterior-mean}

\section{Background}\label{background-1}

This chapter covers the Bayesian v Frequentist methods for estimating
parameters. Imagine you flip a potentially weighted coin N times, and X
are heads. How would you about estimating the true probability that the
coin lands on heads? This is an interesting question, and Bayesians and
frequentists would take different approaches.

This serves as a motivation for estimating parameters. Below, I will
estimate \(p\) in the Negative Binomial Distribution.

First, I will first use the Maximum Likelihood Estimation method to
estimate \(p\), which is the frequentist approach. Then, I will use a
posterior mean, which is the Bayesian approach. Finally, I will compare
them using my dataset.

\section{Maximum Likelihood
Estimation}\label{maximum-likelihood-estimation}

The frequentist method for estimating parameter is finding the MLE.
Below I derive the MLE for \(p\) in the negative binomial.

\(X_1,...,X_n \sim \text{NegBin}(r, p)\)

\[
\begin{aligned}
    P(X=x) &= \binom{x + r- 1}{x} p^r (1 - p)^x \\
    L_{x_1,...,x_n}(p) &= \prod_{i=1}^n \binom{x_i + r- 1}{x_i} p^r (1 - p)^{x_i} \\
    L_{x_1,...,x_n}(p) &= p^{rn} \prod_{i=1}^n \binom{x_i + r- 1}{x_i}  (1 - p)^{x_i} \\
    L_{x_1,...,x_n}(p) &= p^{rn} \Bigg(\prod_{i=1}^n \binom{x_i + r- 1}{x_i} \Bigg) (1 - p)^{\sum_{i=1}^n x_i} \\
    l_{x_1,...,x_n}(p) &= rn \log(p) + \sum_{i=1}^n \log \left(\frac{(x_i + r - 1)!}{{x_i}!(r - 1)!}\right) + \left(\sum_{i=1}^n x_i\right) \log(1 - p) \\
    \frac{d}{dp} l_{x_1,...,x_n}(p) &= \frac{d}{dp} \left( rn\log(p) + \sum_{i=1}^n \log \left(\frac{(x_i + r - 1)!}{{x_i}!(r - 1)!}\right) + \left(\sum_{i=1}^n x_i\right) \log(1 - p) \right) \\
    \frac{d}{dp} l_{x_1,...,x_n}(p) &= \left(\frac{rn}{p} - \frac{\sum_{i=1}^n x_i}{1 - p} \right) \\
    0 &= \frac{rn}{p} - \frac{\sum_{i=1}^n x_i}{1-p} \\
    \frac{rn}{p} &= \frac{\sum_{i=1}^n x_i}{1-p}  \\
    \frac{1-p}{p} &= \frac{\sum_{i=1}^n x_i}{rn} \\
    \frac{1}{p} - 1 &= \frac{\bar{x}}{r} \\
    \hat{p} &= \frac{r}{r + \bar{x}}
\end{aligned}
\]

\section{Posterior Mean}\label{posterior-mean}

For the Bayesian method, we need a prior distribution that we will then
have the data update. Let's use a Beta Prior with
\(\alpha=1, \beta = 1\) (which is actually just a uniform).

\[
\begin{aligned}
    f(p | x) &= \frac{f(p)f(x_1,...,x_n|p)}{f(x_1,...,x_n)}  \\
    f(p | x) &= \frac{\frac{\Gamma(\alpha + \beta)}{\Gamma(\alpha)\Gamma(\beta)} p^{\alpha - 1} (1-p)^{\beta - 1} \prod_{i=1}^n \binom{x_i + r- 1}{x_i} p^r (1 - p)^{x_i}}{f(x_1,...,x_n)} \\ 
    f(p | x) &\propto p^{\alpha + nr - 1} (1-p)^{\beta + \sum x_i - 1}
\end{aligned}
\]

This results in
\(f(p | x) \sim \text{Beta}(\alpha + nr, \beta + \sum_{i=1}^n x_i)\).

\[
\begin{aligned}
    \mathbb{E}[P] &= \frac{\alpha + nr}{\alpha + \beta + nr + \sum_{i=1}^n x_i} \\
    \mathbb{E}[P] &= \frac{1 + nr}{2 + nr + \sum_{i=1}^n x_i}
\end{aligned}
\]

\section{My Dataset}\label{my-dataset}

Now, let's compute the two estimators on my dataset.

\begin{Shaded}
\begin{Highlighting}[]
\NormalTok{dat \textless{}{-} read.csv("https://raw.githubusercontent.com/Ryder4real/STATS200Q\_Book/main/transmission\_pairs\_covid\_v2.csv")}
\NormalTok{x \textless{}{-} as.vector(table(dat$infector\_id)) {-} 1}

\NormalTok{x\_bar  \textless{}{-} mean(x)}
\NormalTok{n      \textless{}{-} length(x)}
\NormalTok{x\_sum  \textless{}{-} sum(x)}

\NormalTok{cat("x\_bar:", x\_bar, "\textbackslash{}n")}
\NormalTok{cat("n:", n, "\textbackslash{}n")}
\NormalTok{cat("x\_sum:", x\_sum, "\textbackslash{}n")}
\end{Highlighting}
\end{Shaded}

As we can see from the code above, in my dataset,
\(\bar{x} = 0.7434944, n=807, \sum(x) = 600\).

So, plugging this in, we have:

\[
\begin{aligned}
    \hat{p}_{MLE} &= \frac{r}{r + 0.7434944} \\
    \hat{p}_{posterior} &= \frac{1 + 807r}{2 + 807r + 600}
\end{aligned}
\]

If we assume \(r = 0.5\):

\[
\begin{aligned}
    \hat{p}_{MLE} &= \frac{0.5}{0.5 + 0.7434944 } = 0.402 \\
    \hat{p}_{posterior} &= \frac{1 + 807(0.5)}{2 + 807(0.5) + 600} = 0.401
\end{aligned}
\]

\section{Comparing Estimators}\label{comparing-estimators}

Now, let's compare the bias and variance of the two estimators. For the
sake of evaluating our estimators, let's say true \(p\) is
\(\frac{1}{3}\).

In my dataset, there are \(807\) patients. After 10,000 simulations, the
bias and variance of our estimators is:

\begin{Shaded}
\begin{Highlighting}[]
\NormalTok{n\_obs \textless{}{-} 807}
\NormalTok{true\_size \textless{}{-} 0.5 \#r}
\NormalTok{true\_p \textless{}{-} 1 / 3}

\NormalTok{sims \textless{}{-} replicate(10000, \{}
\NormalTok{  x \textless{}{-} rnbinom(n = n\_obs, size = true\_size, prob = true\_p)}
\NormalTok{  x\_sum \textless{}{-} sum(x)}
\NormalTok{  x\_avg = mean(x)}
\NormalTok{  p.mle \textless{}{-} true\_size / (true\_size + x\_avg)}
  
\NormalTok{  p.pm \textless{}{-} (1 + n\_obs * true\_size) / (2 + n\_obs * true\_size + x\_sum)}
  
\NormalTok{  c(mle = p.mle, pm = p.pm)}
\NormalTok{\})}

\NormalTok{sims \textless{}{-} t(sims)}

\NormalTok{bias\_mle \textless{}{-} mean(sims[,"mle"]) {-} true\_p}
\NormalTok{bias\_pm \textless{}{-} mean(sims[,"pm"]) {-} true\_p}

\NormalTok{var\_mle \textless{}{-} var(sims[,"mle"])}
\NormalTok{var\_pm \textless{}{-} var(sims[,"pm"])}

\NormalTok{results \textless{}{-} data.frame(}
\NormalTok{  Metric = c("Bias", "Variance"),}
\NormalTok{  MLE = c(bias\_mle, var\_mle),}
\NormalTok{  Posterior\_Mean = c(bias\_pm, var\_pm)}
\NormalTok{)}

\NormalTok{results}
\end{Highlighting}
\end{Shaded}

As we can see from the table produced by the code, \(\hat{p}_{MLE}\) has
positive bias while \(\hat{p}_{PM}\) is biased toward the prior
distribution (which in this case leads to positive bias).

\(\hat{p}_{PM}\) has less slightly less variance. This is more apparent
if you make \texttt{n\_obs} smaller in the code.

Here is the the code for a histogram of \(\hat{p}_{MLE}\) over many
trials.

\begin{Shaded}
\begin{Highlighting}[]
\NormalTok{hist\_mle \textless{}{-} hist(sims[, "mle"], }
\NormalTok{     breaks = 30, }
\NormalTok{     main = expression("Histogram of "* hat(p)[MLE]), }
\NormalTok{     xlab = expression(hat(p)[MLE])}
\NormalTok{    )}
\end{Highlighting}
\end{Shaded}

Here is the the code for a histogram of \(\hat{p}_{PM}\) over many
trials.

\begin{Shaded}
\begin{Highlighting}[]
\NormalTok{hist\_pm \textless{}{-} hist(sims[, "pm"], }
\NormalTok{     breaks = 30, }
\NormalTok{     main = expression("Histogram of "* hat(p)[PM]), }
\NormalTok{     xlab = expression(hat(p)[PM])}
\NormalTok{     )}
\end{Highlighting}
\end{Shaded}

\bookmarksetup{startatroot}

\chapter{Further Evaluating
Estimators}\label{further-evaluating-estimators}

\section{Background}\label{background-2}

This chapter covers a variety of topics. First, I will dive into some
cool properties regarding the score, which is the derivative of the log
likelihood. One integration technique I will use for this in
``Differentiation Under the Integral Sign''. It is a cool trick that
allows you to move a derivative outside of the integral.

I will then dive into another way to evaluate estimators. In Chapter 1,
we looked at the bias and variance of \(\hat{p}_{MLE}\) and
\(\hat{p}_{PM}\), but this time I will use frequentist Mean Squared
Error to compare them. This is defined as MSE =
\(\mathbb{E}[(\hat{\theta} - \theta)^2]\). The MSE can famously be
decomposed as the (Bias\(^2\) + Variance). So, while last chapter we
looked at bias and variance, MSE serves as an all-encompassing metric.
So, I will compare the MSE of my MLE and Posterior Mean from Chapter 1.

Next, I will prove the Cramer-Rao Lower Bound (CRLB), evaluate the
Cramer-Rao Lower Bound (CRLB) for my distribution, and see if it is
relevant to the MLE and Posterior Mean I calculated in Chapter 1. In
class, we learned that the CRLB is a lower bound to the variance of an
unbiased estimator. This means of all unbiased estimators, if one
reaches the CRLB, then it is the best unbiased estimator. This is very
insightful. It means once we've reached that bound, there is no point
evaluating other unbiased estimators, as there cannot be a unbiased
estimator with a lower variance that the one that reaches the CRLB.

As an extra note, the CRLB is defined as \(\frac{1}{I(\theta)}\), where
\(I(\theta)\) is defined as the variance of the score (which is the
derivative of the log-likelihood). I will prove this later on in this
chapter. It is important to note that some of the equations I derive in
the first part of this chapter will make it easier to solve for the CRLB
for my distribution.

\section{\texorpdfstring{Proving
\(\text{Var}(\ell'(\theta)) = -\mathbb{E}[\ell''(\theta)]\)}{Proving \textbackslash text\{Var\}(\textbackslash ell\textquotesingle(\textbackslash theta)) = -\textbackslash mathbb\{E\}{[}\textbackslash ell\textquotesingle\textquotesingle(\textbackslash theta){]}}}\label{proving-textvarelltheta--mathbbeelltheta}

For any PMF/PDF \(f_\theta(x)\), the log-likelihood of \(x\) can be
written as \(\ell_x(\theta)\).

Another name for the derivative of the log-likelihood is the score. The
expected value of the score is defined as
\(\mathbb{E}[\ell'_\theta(x)]\).

\[
\mathbb{E}[\ell'_x(\theta)] = \int_{-\infty}^\infty \ell'_x(\theta)f_\theta(x) dx
\]

Now, let's solve for \(\ell'_x(\theta)\):

\[
\begin{aligned}
L_x(\theta) &= f_\theta(x) \\ 
\ell_x(\theta) &= \log(f_\theta(x)) \\
\frac{d}{d \theta} \ell_x(\theta) &= \frac{\frac{d}{d \theta} f_\theta(x)}{f_\theta(x)} \\
\ell'_x(\theta) &= \frac{\frac{d}{d \theta} f_\theta(x)}{f_\theta(x)}
\end{aligned}
\]

Subbing this into our expectation:

\[
\begin{aligned}
\mathbb{E}[\ell'_x(\theta)] &= \int_{-\infty}^\infty \frac{\frac{d}{d \theta} f_\theta(x)}{f_\theta(x)} f_\theta(x) dx \\
&= \int_{-\infty}^\infty \frac{d}{d \theta} f_\theta(x) dx \\
&= \frac{d}{d \theta} \int_{-\infty}^\infty f_\theta(x) dx \\
&= \frac{d}{d \theta} (1) = 0
\end{aligned}
\] The fact that \(\mathbb{E}[\ell'_x(\theta)] = 0\) is an important
result and I use in again later on in this chapter.

Now, let's solve for \(\text{Var}(\ell'(\theta))\).

We know:
\(\text{Var}(\ell'(\theta)) = \mathbb{E}[(\ell'(\theta))^2] - (\mathbb{E}[\ell'(\theta)])^2\).

Since \(\mathbb{E}[\ell'_x(\theta)] = 0\):

\[
\text{Var}(\ell'(\theta)) = \mathbb{E}[(\ell'(\theta))^2] = \int_{-\infty}^\infty (\ell'_x(\theta))^2 f_\theta(x) dx
\]

Next, we relate \((\ell'_x(\theta))^2\) to the second derivative:

\[
\begin{aligned}
\ell''_x(\theta) = \frac{d}{d \theta} \left[ \frac{\frac{d}{d \theta} f_\theta(x)}{f_\theta(x)} \right] &= \frac{f_\theta(x) \frac{d^2}{d \theta^2} f_\theta(x) - (\frac{d}{d \theta} f_\theta(x))^2}{(f_\theta(x))^2} \\
\ell''_x(\theta) &= \frac{\frac{d^2}{d \theta^2} f_\theta(x)}{f_\theta(x)} - \left( \frac{\frac{d}{d \theta} f_\theta(x)}{f_\theta(x)} \right)^2 \\
\ell''_x(\theta)&= \frac{\frac{d^2}{d \theta^2} f_\theta(x)}{f_\theta(x)} - (\ell'_x(\theta))^2
\end{aligned}
\]

Rearranging gives
\((\ell'_x(\theta))^2 = \frac{\frac{d^2}{d \theta^2} f_\theta(x)}{f_\theta(x)} - \ell''_x(\theta)\).

Subbing this into
\(\text{Var}(\ell'(\theta)) = \mathbb{E}[(\ell'(\theta))^2] = \int_{-\infty}^\infty (\ell'_x(\theta))^2 f_\theta(x) dx\),
gives us:

\[
\begin{aligned}
\text{Var}(\ell'(\theta)) &= \int_{-\infty}^\infty f_\theta(x) \left( \frac{\frac{d^2}{d \theta^2} f_\theta(x)}{f_\theta(x)} - \ell''_x(\theta) \right) dx \\
&= \int_{-\infty}^\infty \frac{d^2}{d \theta^2} f_\theta(x) -  \ell''_x(\theta) f_\theta(x) dx \\
&= \int_{-\infty}^\infty \frac{d^2}{d \theta^2} f_\theta(x)  dx - \int_{-\infty}^\infty \ell''_x(\theta) f_\theta(x) dx \\
&= \frac{d^2}{d\theta^2} \int_{-\infty}^\infty f_\theta(x) dx- \mathbb{E}[\ell''_x(\theta)] \\
&= \frac{d^2}{d\theta^2} (1) - \mathbb{E}[\ell''_x(\theta)] \\
&= - \mathbb{E}[\ell''_x(\theta)]
\end{aligned}
\]

So, we have just shown that
\(\text{Var}(\ell'(\theta)) = - \mathbb{E}[\ell''_x(\theta)]\). This
value is also known as the Fisher Information, or \(I(\theta)\).

\section{MSE}\label{mse}

In Chapter 1, I derived that if \(X \sim \text{NegBin}(r, p)\), then:
\(\hat{p}_{MLE} = \frac{r}{r + \bar{x}}\) and \(\hat{p}_{PM}\) =
\(\frac{1 + nr}{2 + nr + \sum x_i}\).

To show that \(\hat{p}_{MLE}\) and \(\hat{p}_{PM}\) have different MSE
for different \(p\), below is a simulation that prints out Mean Squared
Error for different \(p\), for when \(n=50\) and \(r=0.5\).

\begin{Shaded}
\begin{Highlighting}[]
\NormalTok{n }\OtherTok{\textless{}{-}} \DecValTok{50} 
\NormalTok{true\_size }\OtherTok{\textless{}{-}} \FloatTok{0.5} \CommentTok{\#r}

\NormalTok{p\_grid }\OtherTok{\textless{}{-}} \FunctionTok{seq}\NormalTok{(}\FloatTok{0.01}\NormalTok{, }\FloatTok{0.99}\NormalTok{, }\AttributeTok{length.out =} \DecValTok{40}\NormalTok{)}
\NormalTok{mse\_results }\OtherTok{\textless{}{-}} \FunctionTok{sapply}\NormalTok{(p\_grid, }\ControlFlowTok{function}\NormalTok{(p) \{}
  
  \CommentTok{\# Run simulations for each specific p in the grid}
\NormalTok{  temp\_sims }\OtherTok{\textless{}{-}} \FunctionTok{replicate}\NormalTok{(}\DecValTok{10000}\NormalTok{, \{}
\NormalTok{    x }\OtherTok{\textless{}{-}} \FunctionTok{rnbinom}\NormalTok{(}\AttributeTok{n =}\NormalTok{ n, }\AttributeTok{size =}\NormalTok{ true\_size, }\AttributeTok{prob =}\NormalTok{ p)}
\NormalTok{    x\_avg }\OtherTok{=} \FunctionTok{mean}\NormalTok{(x)}
\NormalTok{    x\_sum }\OtherTok{=} \FunctionTok{sum}\NormalTok{(x)}
    
\NormalTok{    p.mle }\OtherTok{\textless{}{-}}\NormalTok{ true\_size }\SpecialCharTok{/}\NormalTok{ (true\_size }\SpecialCharTok{+}\NormalTok{ x\_avg)}
\NormalTok{    p.pm }\OtherTok{\textless{}{-}}\NormalTok{ (}\DecValTok{1} \SpecialCharTok{+}\NormalTok{ n }\SpecialCharTok{*}\NormalTok{ true\_size) }\SpecialCharTok{/}\NormalTok{ (}\DecValTok{2} \SpecialCharTok{+}\NormalTok{ n }\SpecialCharTok{*}\NormalTok{ true\_size }\SpecialCharTok{+}\NormalTok{ x\_sum)}
    
    \FunctionTok{c}\NormalTok{(}\AttributeTok{mle =}\NormalTok{ p.mle, }\AttributeTok{pm =}\NormalTok{ p.pm)}
\NormalTok{  \})}
  
  \CommentTok{\# Calculate MSE for this specific p}
\NormalTok{  mse\_mle\_p }\OtherTok{\textless{}{-}} \FunctionTok{mean}\NormalTok{((temp\_sims[}\StringTok{"mle"}\NormalTok{,] }\SpecialCharTok{{-}}\NormalTok{ p)}\SpecialCharTok{\^{}}\DecValTok{2}\NormalTok{)}
\NormalTok{  mse\_pm\_p }\OtherTok{\textless{}{-}} \FunctionTok{mean}\NormalTok{((temp\_sims[}\StringTok{"pm"}\NormalTok{,] }\SpecialCharTok{{-}}\NormalTok{ p)}\SpecialCharTok{\^{}}\DecValTok{2}\NormalTok{)}
  
  \FunctionTok{c}\NormalTok{(}\AttributeTok{mse\_mle =}\NormalTok{ mse\_mle\_p, }\AttributeTok{mse\_mod =}\NormalTok{ mse\_pm\_p)}
\NormalTok{\})}

\CommentTok{\# Plotting the curve}
\FunctionTok{plot}\NormalTok{(p\_grid, mse\_results[}\StringTok{"mse\_mle"}\NormalTok{,], }\AttributeTok{type =} \StringTok{"l"}\NormalTok{, }\AttributeTok{col =} \StringTok{"blue"}\NormalTok{, }\AttributeTok{lwd =} \DecValTok{2}\NormalTok{,}
     \AttributeTok{main =} \StringTok{"MSE for Different Estimators"}\NormalTok{, }
     \AttributeTok{xlab =} \StringTok{"True p"}\NormalTok{, }
     \AttributeTok{ylab =} \StringTok{"MSE"}\NormalTok{,}
     \AttributeTok{cex.main =} \FloatTok{0.9}\NormalTok{,  }
     \AttributeTok{cex.lab =} \FloatTok{0.8}\NormalTok{,   }
     \AttributeTok{cex.axis =} \FloatTok{0.7}\NormalTok{)}

\FunctionTok{lines}\NormalTok{(p\_grid, mse\_results[}\StringTok{"mse\_mod"}\NormalTok{,], }\AttributeTok{col =} \StringTok{"red"}\NormalTok{, }\AttributeTok{lwd =} \DecValTok{2}\NormalTok{, }\AttributeTok{lty =} \DecValTok{2}\NormalTok{)}
\FunctionTok{legend}\NormalTok{(}\StringTok{"topright"}\NormalTok{, }
       \AttributeTok{legend =} \FunctionTok{c}\NormalTok{(}\StringTok{"MLE"}\NormalTok{, }\StringTok{"Posterior Mean"}\NormalTok{), }
       \AttributeTok{col =} \FunctionTok{c}\NormalTok{(}\StringTok{"blue"}\NormalTok{, }\StringTok{"red"}\NormalTok{), }
       \AttributeTok{lty =} \DecValTok{1}\SpecialCharTok{:}\DecValTok{2}\NormalTok{,}
       \AttributeTok{cex =} \FloatTok{0.6}\NormalTok{)}
\end{Highlighting}
\end{Shaded}

\pandocbounded{\includegraphics[keepaspectratio]{Week2_files/figure-pdf/unnamed-chunk-1-1.pdf}}

As we can see, at different true \(p\) values, one estimator may be
better than the other.

\section{Proving the Cramer-Rao Lower
Bound}\label{proving-the-cramer-rao-lower-bound}

As mentioned in the ``Background'' section, the CRLB is a lower bound to
the variance of an unbiased estimator. Let's show why this is the case.

First, let \(t(x)\) be any unbiased estimator of \(\theta\). Now let's
solve for \(\text{Cov}(t(x), \ell'_x(\theta))\).

\[
\begin{aligned}
\text{Cov}(t(x), \ell'(x)) &= \mathbb{E}[t(x) \cdot \ell'_x(\theta)] - \mathbb{E}[t(x)] \mathbb{E}[\ell'_x(\theta)] \\
&= \mathbb{E}[t(x) \cdot \ell'_x(\theta)] - \mathbb{E}[t(x)] \mathbb{E}[\ell'_x(\theta)]
\end{aligned}
\] Now, in the first section of this Chapter, we showed that
\(\mathbb{E}[\ell'_x(\theta)] = 0\). So, let's apply that and continue
to manipulate.

\[
\begin{aligned}
\text{Cov}(t(x), \ell'(x)) &= \mathbb{E}[t(x) \cdot \ell'_x(\theta)] - 0 \\
&= \int_{-\infty}^\infty t(x) \ell'(x) f_\theta(x) dx
\end{aligned}
\]

Now, in the first section we also showed that:
\(\ell'_x(\theta) = \frac{\frac{d}{d \theta} f_\theta(x)}{f_\theta(x)}\)

Using this, we can say that:

\[
\begin{aligned}
\text{Cov}(t(x), \ell'(x)) &= \int_{-\infty}^\infty t(x) \frac{\frac{d}{d \theta} f_\theta(x)}{f_\theta(x)} f_\theta(x) dx \\
&= \int_{-\infty}^\infty t(x) \frac{d}{d \theta} f_\theta(x) dx \\
&= \frac{d}{d \theta}\int_{-\infty}^\infty t(x)  f_\theta(x) dx \\
&= \frac{d}{d \theta}\mathbb{E}[t(x)] \\
&= \frac{d}{d \theta} \theta \text{ [NOTE: This is because t(x) is unbiased]} \\
&= 1
\end{aligned}
\] So, using the Cauchy-Schwarz Inequality, we can say that

\(1 = \text{Cov}^2(t(x), \ell'(x)) \leq \text{Var}(t(x)) \text{Var}(\ell'_x(\theta))\)

So, this means that:

\(\frac{1}{\text{Var}(\ell'_x(\theta))} \leq \text{Var}(t(x))\).

So, this means that the variance of any unbiased estimator \(t(x)\) must
be greater than or equal to \(\frac{1}{\text{Var}(\ell'_x(\theta))}\).
And, if you remember earlier, this can be written as:

CRLB:
\(\frac{1}{\text{Var}(\ell'_x(\theta))} = \frac{1}{I(\theta)} = \frac{1}{-\mathbb{E}[\ell''_x(\theta)]}\)

\section{Finding the CRLB}\label{finding-the-crlb}

As described above, the CRLB for unbiased estimators is
\(\frac{1}{I(\theta)}\), where \(I(\theta)\) is defined as the Variance
of the Score. However, we proved above that the variance of the score,
\(\text{Var}(\ell'(\theta))\) is actually equivalent to
\(-\mathbb{E}[\ell''(\theta)]\)

So, this means that the CRLB for unbiased estimators is
\(\frac{1}{-\mathbb{E}[\ell''(\theta)]}\).

In Chapter 1, I wrote For \(X_1,...,X_n \sim \text{NegBin}(r, p)\), then
\(\ell'(p) = \frac{rn}{p} - \frac{\sum_{i=1}^n x_i}{1-p}\):

This means that
\(\ell''(p) = -\frac{rn}{p^2} - \frac{\sum_{i=1}^n x_i}{(1-p)^2}\)

\[
\begin{aligned}
\mathbb{E}\left[\ell''(p)\right] &= -\left(\frac{rn}{p^2} + \frac{\mathbb{E}[\sum_{i=1}^n x_i]}{(1-p)^2} \right) \\
&= -\left(\frac{rn}{p^2} + \frac{rn(1-p)}{p(1-p)^2} \right) \\
&= -\frac{rn}{p^2(1-p)}
\end{aligned}
\] So, the lowest possible variance for an unbiased estimator for \(p\)
is \(\frac{p^2(1-p)}{rn}\).

\section{Does it Apply?}\label{does-it-apply}

In Week 1, we established from simulation that \(\hat{p}_{MLE}\) is
positively biased. Additionally, in class we learned that the
\(\hat{p}_{PM}\) is always biased toward the prior. Since the MLE and
Posterior Mean are both biased in this case, the CRLB does not provide a
strict lower bound for their variance.

\bookmarksetup{startatroot}

\chapter{Asymptotic Distribution of the
MLE}\label{asymptotic-distribution-of-the-mle}

\section{Background}\label{background-3}

This chapter will dive into the Asymptotic Distribution of the MLE we
calculated in Chapter 1. This is interesting because calculating the
bias and variance of the MLE by hand can be difficult.

However, this can be made easier when dealing with large sample sizes by
knowing that the distribution of the MLE approximates to a Normal as
\(n \rightarrow \infty\) It is kind of wild to think about no matter the
original distribution, the distribution of the MLE always approaches the
Normal distribution!

I will first simulate my distribution and look at a histogram of
\(\hat{p}_{MLE}\) as \(n\) grows large. Then, I will prove that the MLE
approximates to a Normal as \(n \rightarrow \infty\) and apply that to
my distribution.

\section{Simulating Distribution}\label{simulating-distribution}

Let's say \(X_1, ...,X_n \sim \text{NegBin}(r, p)\), where \(r=0.5\) and
\(p = \frac{1}{3}\).

We derived in Chapter 1. that \(\hat{p}_{MLE} = \frac{r}{r + \bar{x}}\).
Below is a histogram of \(\hat{p}_{MLE}\) as \(n\) grows large, in this
case to \(10,000\).

\begin{Shaded}
\begin{Highlighting}[]
\NormalTok{n\_obs \textless{}{-} 10000  }
\NormalTok{true\_size \textless{}{-} 0.5 \#r}
\NormalTok{true\_p \textless{}{-} 1 / 3}

\NormalTok{sims \textless{}{-} replicate(1000, \{}
\NormalTok{  x \textless{}{-} rnbinom(n = n\_obs, size = true\_size, prob = true\_p)}
\NormalTok{  x\_sum \textless{}{-} sum(x)}
\NormalTok{  x\_avg = mean(x)}
\NormalTok{  p.mle \textless{}{-} true\_size / (true\_size + x\_avg)}
  
\NormalTok{  p.mle}
\NormalTok{\})}

\NormalTok{sims \textless{}{-} t(sims)}

\NormalTok{hist\_mle \textless{}{-} hist(sims, }
\NormalTok{     breaks = 100, }
\NormalTok{     main = expression("Histogram of "* hat(p)[MLE]), }
\NormalTok{     xlab = expression(hat(p)[MLE])}
\NormalTok{)}
\end{Highlighting}
\end{Shaded}

As we can see, this looks pretty Normal, and it looks to be centered
around \(\frac{1}{3}\), which is what we set true \(p\) to be.

\section{Proof of MLE Approaching a Normal
Distribution}\label{proof-of-mle-approaching-a-normal-distribution}

Now let's prove that the distribution of the MLE approaches a Normal
distribution as \(n\) gets large. This proof is a really cool one and
involves many different concepts.

Let's start with the score of \(\hat{\theta}_{MLE}\), which is
\(\ell'(\hat{\theta}_{MLE})\)

Let's now Taylor expand \(\ell'(\hat{\theta}_{MLE})\) around
\(\theta_0\).

\(\ell'(\hat{\theta}_{MLE}) \approx \ell'(\theta_0) + \ell''(\theta_0)(\hat{\theta}_{MLE} - \theta_0)\)

Now, remember that when solving for the MLE (as we did in Chapter 1), we
set \(\ell'(x) = 0\). This is because the MLE is defined as the
\(\theta\) value that maximizes the likelihood (and log-likelihood).

So, as the MLE is a maximum of the log-likelihood function, this means
that \(\ell'(\hat{\theta}_{MLE}) = 0\).

So, plugging this in we get:

\(0 = \ell'(\hat{\theta}_{MLE}) \approx \ell'(\theta_0) + \ell''(\theta_0)(\hat{\theta}_{MLE} - \theta_0)\)

We can rearrange this to say:

\[
\begin{aligned}
  0 &\approx \ell'(\theta_0) + \ell''(\theta_0)(\hat{\theta}_{MLE} - \theta_0) \\
  - \ell'(\theta_0) &\approx \ell''(\theta_0)(\hat{\theta}_{MLE} - \theta_0) \\
  \hat{\theta}_{MLE} - \theta_0  &\approx \frac{-\ell'(\theta_0)}{\ell''(\theta_0)}
\end{aligned}
\]

Now, let's do some further manipulation.

\[
\begin{aligned}
  \hat{\theta}_{MLE} - \theta_0  &\approx \frac{-\ell'(\theta_0)}{\ell''(\theta_0)} \\
    \hat{\theta}_{MLE} - \theta_0  &= \frac{\frac{1}{n}\ell'(\theta_0)}{-\frac{1}{n}\ell''(\theta_0)} \\ 
      \hat{\theta}_{MLE} - \theta_0  &=  \frac{\frac{1}{n}\sum_{i=1}^n \ell'_{x_i}(\theta_0)}{-\frac{1}{n}\sum_{i=1}^n\ell''_{x_i}(\theta_0)} \\ 
      \sqrt{n}(\hat{\theta}_{MLE} - \theta_0)  &= \frac{\sqrt{n}(\frac{1}{n}\sum_{i=1}^n \ell'_{x_i}(\theta_0))}{-\frac{1}{n}\sum_{i=1}^n\ell''_{x_i}(\theta_0)} \\ 
\end{aligned}
\]

Now, let's say \(Y = \ell'_{x_1}(\theta_0)\). This means that
\(\bar{Y} = \frac{1}{n} \sum_{i=1}^n \ell'_{x_i}(\theta_0)\).

This means that:

\[
\begin{aligned}
  \sqrt{n}(\hat{\theta}_{MLE} - \theta_0)  &= \frac{\sqrt{n} (\bar{Y} - 0)} {-\frac{1}{n}\sum_{i=1}^n\ell''_{x_i}(\theta_0)} \\ 
\end{aligned}
\]

Now, we can realize three things.

First, the CLT is defined as
\(\sqrt{n}(\bar{Y} - \mathbb{E}[Y_1]) \rightarrow N(0, Var(Y_1))\)
\textbackslash{} Second, as we showed in Chapter 2,
\(\mathbb{E}[Y_1] = \mathbb{E}[\ell'_{x_1}(\theta_0)] = 0\).
\textbackslash{} Third, also from Chapter 2, we showed that
\(I(\theta) = -\mathbb{E}[\ell''(\theta)] = \text{Var}(\ell'(\theta))\)

So, this means that:

\[
\begin{aligned}
   \sqrt{n}(\hat{\theta}_{MLE} - \theta_0) \rightarrow \frac{N(0, \text{Var}(\ell'_{x_1}(\theta_0))}{\text{Var}(\ell'(\theta_0))} \\ 
    \sqrt{n}(\hat{\theta}_{MLE} - \theta_0) \rightarrow N \Big(0, \frac{1}{\text{Var}(\ell'_{x_1}(\theta_0))}\Big) \\ 
  \sqrt{n}(\hat{\theta}_{MLE} - \theta_0) \rightarrow N \Big(0, \frac{1}{I_{x_1}(\theta)}\Big) \\ 
\end{aligned}
\]

Isolating for \(\hat{p}_{MLE}\), we have that
\(\hat{p}_{MLE} \rightarrow N\left(p, \frac{1}{nI_{x_1}(\theta)}\right)\)

\section{Asymptotic Calculation}\label{asymptotic-calculation}

As shown above, using the CLT and Taylor approximation, we can see that
as \(n\) grows large, we have that
\(\hat{p}_{MLE} \rightarrow N\left(p, \frac{1}{nI_{x_1}(\theta))}right)\)

So, let's solve for \(I_{x_1}(\theta)\). Once again, from Chapter 2, we
know that
\(I(\theta) = -\mathbb{E}[\ell''(\theta)] = \text{Var}(\ell'(\theta))\)

Taking the score that we got from Chapter 1 (while solving for the MLE),
we have:

\[
\frac{d}{dp} \ell_{x_1,...,x_n}(p) = \left(\frac{rn}{p} - \frac{\sum x}{1 - p} \right)
\]

As we only want to work with one observation now, we will set \(n=1\).

\[
\ell ^\prime_{x_1} (p) = \left(\frac{r}{p} - \frac{x}{1 - p} \right)
\]

Let's solve for \(-\mathbb{E}[\ell''(\theta)]\).

\[
\begin{aligned}
\ell ^\prime_{x_1} (p) &= \frac{r}{p} - \frac{x}{1 - p}  \\
\frac{d^2}{dp^2} \ell_{x_1} (p) &= \frac{-r}{p^2} - \frac{x}{(1 - p)^2} \\
-\mathbb{E}[\ell ^{\prime \prime}_{x_1} (p)] &= \frac{r}{p^2} + \frac{\mathbb{E}[x]}{(1 - p)^2} \\
-\mathbb{E}[\ell ^{\prime \prime}_{x_1} (p)] &= \frac{r}{p^2} + \frac{\Big(\frac{r(1-p)}{p}\Big)}{(1 - p)^2} \\
-\mathbb{E}[\ell ^{\prime \prime}_{x_1} (p)] &= \frac{r}{p^2} + \frac{r}{p(1 - p)} \\
-\mathbb{E}[\ell ^{\prime \prime}_{x_1} (p)] &= \frac{r(1-p)}{p^2(1-p)} + \frac{rp}{p^2(1 - p)} \\
-\mathbb{E}[\ell ^{\prime \prime}_{x_1} (p)] &= \frac{r(1-p) + rp}{p^2(1-p)}  \\
-\mathbb{E}[\ell ^{\prime \prime}_{x_1} (p)] &= \frac{r(1 - p + p)}{p^2(1-p)}  \\
-\mathbb{E}[\ell ^{\prime \prime}_{x_1} (p)] &= \frac{r}{p^2(1-p)}  \\
\end{aligned}
\]

So, we have

\[
  \sqrt{n} (\hat{p}_{MLE} - p) \rightarrow N\left(0, \frac{p^2(1-p)}{r}\right)
\]

And again isolating for \(\hat{p}_{MLE}\), we have:

\[
\hat{p}_{MLE} \overset{\cdot}{\sim} N\left(p, \frac{p^2(1-p)}{nr} \right)
\]

\section{Graphing Asymptotic MLE
Curve}\label{graphing-asymptotic-mle-curve}

Below is a curve with the distribution of \(\hat{p}_{MLE}\) with
\(r=0.5\), \(p= \frac{1}{3}\), and large \(n (10,000)\), behind a red
curve displaying the asymptotically normal distribution of
\(\hat{p}_{MLE}\).

\begin{Shaded}
\begin{Highlighting}[]
\NormalTok{hist\_mle \textless{}{-} hist(sims, }
\NormalTok{     breaks = 100, }
\NormalTok{     main = expression("Histogram of "* hat(p)[MLE]), }
\NormalTok{     xlab = expression(hat(p)[MLE]),}
\NormalTok{     freq=F}
\NormalTok{)}
\NormalTok{curve(dnorm(x, mean=true\_p,}
\NormalTok{            sd= sqrt( ( (true\_p ** 2) * (1 {-} true\_p) ) / (n\_obs * true\_size))),}
\NormalTok{add=T, col="red", lwd=2)}

\NormalTok{legend("topright", legend = expression("Asymptotic Normal "* hat(p)[MLE]), }
\NormalTok{       col="red", lwd=2, lty = 1, cex = 0.6)}
\end{Highlighting}
\end{Shaded}

\bookmarksetup{startatroot}

\chapter{Confidence and Credible
Intervals}\label{confidence-and-credible-intervals}

\section{Background}\label{background-4}

So far, we have discussed solving for and evaluating \(\hat{p}_{MLE}\)
and \(\hat{p}_{PM}\). But, these estimators are based only on the data w
ehave, and neither is likely to be the true \(p\) of the entire
population.

So, this chapter is all about Confidence and Credible Intervals. The
former is a frequentist method to display a region that they believe
contains the true \(p\) of the entire population with a certain
confidence. The latter is the Bayesian version that allows for slightly
different wording that I will explain in the next paragraph. For these
examples, we use a \(95\%\) confidence and credible interval.

This is the topic that truly instilled in me the differences in
frequentist v. Bayesian principles. As frequentists treat \(p\) as a
fixed but unknown parameter, once they have calculated their region,
they do NOT say that there is a \(95\%\) chance that \(p\) is in the
region. This is because \(p\) is fixed. A fixed parameter cannot have a
chance of being in a region, it either is or it isn't. So, they escape
this by saying they are \(95\%\) ``confident'' that \(p\) is contained
in the region. On the the other hand, as Bayesians treat \(p\) is a
random variable, they have no problems saying that there is a \(95\%\)
chance that \(p\) is in their credible interval.

\section{Confidence Interval}\label{confidence-interval}

To remember, in Chapter 1, we solved that for
\(X_1,...,X_n \sim \text{NegBin}(r, p)\)
\(\hat{p}_{MLE} = \frac{r}{r + \bar{x}}\).

In Chapter 3, we established both that

\[
  \sqrt{n} (\hat{p}_{MLE} - p) \sim N\left(0, \frac{p^2(1-p)}{r}\right)
\]

and

\[
\hat{p}_{MLE} \overset{\cdot}{\sim} N\left(p, \frac{p^2(1-p)}{nr} \right)
\]

This means that as the number of data points approaches infinity,
\(\hat{p}_{MLE}\) is approximately Normal, with a mean of \(p\) and a
variance of \(\frac{p^2(1-p)}{nr}\).

Therefore, an approximate 95\% confidence interval for \(p\) is:
\(\hat{p}_{MLE} \pm 2 \sqrt{\frac{p^2(1-p)}{nr}}\). However, we do not
know the true value of \(p\), so we have two options in how we can
proceed.

\subsection{Plug-In}\label{plug-in}

We can try plugging in \(\hat{p}_{MLE}\) for \(p\) in this equation.
This means that the approximate 95\% confidence interval is:

\[
\begin{aligned}
\hat{p} \pm 2 \sqrt{\frac{\hat{p}_{MLE}^2(1-\hat{p}_{MLE})}{nr}} 
&= \hat{p} \pm 2 \sqrt{\frac{\left( \frac{r}{r + \bar{x}} \right)^2 \left( 1-\frac{r}{r + \bar{x}} \right)}{nr}} \\
&= \hat{p} \pm 2 \sqrt{\frac{\left( \frac{r}{r + \bar{x}} \right)^2 \left( \frac{r + \bar{x}}{r + \bar{x}}-\frac{r}{r + \bar{x}} \right)}{nr}} \\
&= \hat{p} \pm 2 \sqrt{\frac{\left( \frac{r}{r + \bar{x}} \right)^2 \left( \frac{\bar{x}}{r + \bar{x}} \right)}{nr}} \\
&= \hat{p} \pm 2 \sqrt{\frac{r^2\bar{x}}{nr(r + \bar{x})^3}} \\
&= \hat{p} \pm 2 \sqrt{\frac{r\bar{x}}{n(r + \bar{x})^3}}
\end{aligned}
\]

Going back to my hospital dataset:

\begin{Shaded}
\begin{Highlighting}[]
\NormalTok{dat \textless{}{-} read.csv("https://raw.githubusercontent.com/Ryder4real/STATS200Q\_Book/main/transmission\_pairs\_covid\_v2.csv")}
\NormalTok{x \textless{}{-} as.vector(table(dat$infector\_id)) {-} 1}
\NormalTok{x\_bar = mean(x)}
\NormalTok{n = length(x)}
\NormalTok{cat("x\_bar =", x\_bar, "\textbackslash{}tn =", n)}
\end{Highlighting}
\end{Shaded}

As we can see, I had \(\bar{x} \approx 0.7435, n=807\). Let's say
\(r = 0.5\).

I then have:

\[
\begin{aligned}
\hat{p} \pm 2 \sqrt{\frac{r\bar{x}}{n(r + \bar{x})^3}} \\
\frac{r}{r+\bar{x}} \pm 2 \sqrt{\frac{r\bar{x}}{n(r + \bar{x})^3}} \\
\frac{0.5}{0.5+0.7435} \pm 2 \sqrt{\frac{(0.5)(0.7435)}{807(0.5 + 0.7435)^3}} \\
0.4021 \pm 2 \sqrt{0.00023958}
\end{aligned}
\]

So, the approximate plug-in 95\% confidence interval for \(p\) is
\([0.3711, 0.4331]\).

\subsection{Direct Computation}\label{direct-computation}

Instead of plugging in \(\hat{p}_{MLE}\) for \(p\), we can try to solve
for what values of \(p\) make up the confidence interval.

We established earlier that:

\(\sqrt{n} (\hat{p}_{MLE} - p) \sim N\left(0, \frac{p^2(1-p)}{r}\right)\)

This means that:
\(\hat{p}_{MLE} - p \sim N\left(0, \frac{p^2(1-p)}{nr}\right)\).

Furthermore, we can say that:

\(\frac{\hat{p}_{MLE} - p}{\sqrt{ \frac{p^2(1-p)}{nr}}} \sim N(0, 1)\)

In the Normal distribution, \(95\%\) of the mass is in the mean
\(\pm 1.96\) Standard deviations. As the Normal above has a standard
deviation of 1, we can set up a 95\% confidence interval that rounds to:

\(-2 \leq \frac{\hat{p}_{MLE} - p}{\sqrt{ \frac{p^2(1-p)}{nr}}} \leq 2\)

\[
\begin{aligned}
\Bigg( \frac{\hat{p}_{MLE} - p}{\sqrt{ \frac{p^2(1-p)}{nr}}} \Bigg)^2 &\leq 4 \\
\frac{(\hat{p}_{MLE} - p)^2}{ \frac{p^2(1-p)}{nr}} &\leq 4 \\
\frac{nr(\hat{p}_{MLE} - p)^2}{p^2(1-p)} &\leq 4 \\
\hat{p}_{MLE}^2 -2\hat{p}_{MLE}p + p^2 &\leq \frac{4}{nr}p^2(1-p) \\
\hat{p}_{MLE}^2 -2\hat{p}_{MLE}p + p^2- \frac{4p^2}{nr} + \frac{4p^3}{nr} &\leq 0 \\
\frac{4}{nr}p^3 + (1- \frac{4}{nr})p^2 - 2\hat{p}_{MLE}p + \hat{p}_{MLE}^2&\leq 0 
\end{aligned}
\]

This is unfortunately not cleanly solvable for \(p\), but we can graph
it. If we say that
\(f(p) = \frac{4}{nr}p^3 + (1- \frac{4}{nr})p^2 - 2\hat{p}_{MLE}p + \hat{p}_{MLE}^2\),
we can look at the zeros to see the values of \(p\) at the edge of the
confidence intervals.

For this sample, let's go back to the dataset I started with, where
\(\bar{x} = 0.7435, n=807\). Let's fix \(r = 0.5\).

\begin{Shaded}
\begin{Highlighting}[]
\NormalTok{r \textless{}{-} 0.5      }

\NormalTok{p\_hat \textless{}{-} r / (r + x\_bar)}

\NormalTok{\# f(p)}
\NormalTok{cubic\_f \textless{}{-} function(p) \{}
\NormalTok{  a \textless{}{-} 4 / (n * r)}
\NormalTok{  (a * p\^{}3) + (1 {-} a) * p\^{}2 {-} (2 * p\_hat) * p + p\_hat\^{}2}
\NormalTok{\}}

\NormalTok{\# roots of f(p) the windows [0, p\_hat] and [p\_hat, 1]}
\NormalTok{lower\_bound \textless{}{-} uniroot(cubic\_f, interval = c(0, p\_hat))$root}
\NormalTok{upper\_bound \textless{}{-} uniroot(cubic\_f, interval = c(p\_hat, 1))$root}


\NormalTok{cat("95\% Confidence Interval: [", round(lower\_bound, 4), ",", round(upper\_bound, 4), "]\textbackslash{}n")}
\end{Highlighting}
\end{Shaded}

From the computation in the code above, we are approximately 95\%
confident that \(p\) is in \([0.3727, 0.4346]\)

\section{Credible Interval}\label{credible-interval}

In Chapter 1, I solved that given a Beta\((\alpha, \beta)\) prior, the
posterior distribution for \(p\) is:

\(f(p | x_1,..,x_n) \sim \text{Beta}(\alpha + nr, \beta + \sum_{i=1}^n x_i)\).

In Chapter 1, I used a uniform prior, meaning \(\alpha =1, \beta=1\).
This means the posterior distribution given a uniform prior is:

\(f(p | x_1,..,x_n) \sim \text{Beta}(1 + nr, 1 + \sum_{i=1}^n x_i)\).

For our 95\% credible interval, we want the interval to cover 95\% of
the distribution with symmetric tails.

So, we need a left bound (LB) and right bound (RB), such that:

\(\int_{0}^{LB} \frac{p^{(1+nr)-1} (1-p)^{(1+\sum x_i)-1}}{B(1+nr, 1+\sum x_i)} dp = 0.025\)

\(\int_{RB}^{1} \frac{p^{(1+nr)-1} (1-p)^{(1+\sum x_i)-1}}{B(1+nr, 1+\sum x_i)} dp = 0.025\)

Unfortunately, we can't solve for LB and RB by hand. However, the
computer can solve it for us, using the \texttt{qbeta} function.

Again, we will use the same values from my hospital dataset, as before.

\begin{Shaded}
\begin{Highlighting}[]
\NormalTok{r \textless{}{-} 0.5}
\NormalTok{sum\_xi \textless{}{-} x\_bar * n }
\NormalTok{credible\_range \textless{}{-} 0.95}

\NormalTok{\# Posterior Parameters based on f(p | x) \textasciitilde{} Beta(1 + nr, 1 + sum x\_i)}
\NormalTok{alpha\_post \textless{}{-} 1 + (n * r)}
\NormalTok{beta\_post  \textless{}{-} 1 + sum\_xi}

\NormalTok{\# Calculate Tail Areas}
\NormalTok{alpha\_level \textless{}{-} 1 {-} credible\_range}
\NormalTok{tail\_area   \textless{}{-} alpha\_level / 2}

\NormalTok{\# Solve for Credible Interval Bounds}
\NormalTok{LB \textless{}{-} qbeta(tail\_area, }
\NormalTok{            shape1 = alpha\_post, }
\NormalTok{            shape2 = beta\_post, }
\NormalTok{            lower.tail = TRUE)}

\NormalTok{RB \textless{}{-} qbeta(tail\_area, }
\NormalTok{            shape1 = alpha\_post, }
\NormalTok{            shape2 = beta\_post, }
\NormalTok{            lower.tail = FALSE)}

\NormalTok{credible\_interval \textless{}{-} c(LB, RB)}
\NormalTok{names(credible\_interval) \textless{}{-} c("Left Bound", "Right Bound")}

\NormalTok{print(credible\_interval)}
\end{Highlighting}
\end{Shaded}

From the computation in the code above, conditional on my uniform prior,
the credible interval states that there is a 95\% chance \(p\) is in
\([0.3722, 0.4328]\).

\bookmarksetup{startatroot}

\chapter{Simulating Confidence
Intervals}\label{simulating-confidence-intervals}

\section{Background}\label{background-5}

In the prevous chapter, we talked about using confidence and credible
intervals. For the frequentist calculation, we assumed large \(n\) and
used a Normal approximation for \(\hat{p}_{MLE}\). For our credible
interval, we chose a conjugate prior so we could solve for the posterior
distribution.

However, what if we didn't want to resort to using a Normal
approximatation for our confidence interval? Additionally, what if we
don't know the our posterior distribution for the credible interval?

Statisticians have had come up with computationally-intense mechanisms
to work around these issues. For frequentists, this involves the
parametric and nonparametric bootstrap. In the former, you assume your
population follows a certain distribution, then use the MLE from your
sample to generate more datasets. In the latter, you re-sample more
datasets from your original sample with replacement. For both, you then
look at the MLEs of your many generated datasets.

The Bayesian Process involves a Metropolis-Hastings Markov Chain Monte
Carlo. Let's say we want to come up with estimator for \(\theta\) in
Distribution B. To do this, we first pick a starting Distribution A and
continuously sample from it. After each sample, we establish a
probability based on the likelihood that they can be generated by
Distributions A \& B. We then keep each sample with this probability.
The idea, is after a burn-in period, despite being generated from
Distribution A, the kept samples will actually resemble Distribution B.
I will explain the intricacies in more detail later in the chapter.

\section{Bootstrapping}\label{bootstrapping}

First, let's bring my hospital dataset into R.

\begin{Shaded}
\begin{Highlighting}[]
\NormalTok{dat \textless{}{-} read.csv("https://raw.githubusercontent.com/Ryder4real/STATS200Q\_Book/main/transmission\_pairs\_covid\_v2.csv")}
\NormalTok{x \textless{}{-} as.vector(table(dat$infector\_id)) {-} 1}
\end{Highlighting}
\end{Shaded}

\subsection{Parametric Bootstrap}\label{parametric-bootstrap}

If we say \(X \sim \text{NegBin}(r, p)\), and we have \(r = 0.5\), then
as I calculated in Chapter 1,
\(\hat{p}_{MLE} = \frac{r}{r + \bar{x}} = \frac{0.5}{0.5 + 0.7435} = 0.402\).

\begin{Shaded}
\begin{Highlighting}[]
\NormalTok{r = 0.5}
\NormalTok{n = length(x)}
\NormalTok{x\_bar = sum(x) / n }
\NormalTok{p\_mle = r / (r + x\_bar)}
\end{Highlighting}
\end{Shaded}

So, using a parametric bootstrap, we can simulate many new datasets from
the Negabtive Binomial with my \(\hat{p}_{MLE}\). Then, we can calculate
the MLE of each of these datasets. A histogram of the MLE values of
these datasets can be seen in the code block below.

\begin{Shaded}
\begin{Highlighting}[]
\NormalTok{p.pboot \textless{}{-} replicate(10000, \{}
\NormalTok{    \# Generate dataset}
\NormalTok{    x.pboot = rnbinom(n, r, p\_mle)}
    
\NormalTok{    \# Calculate new MLE}
\NormalTok{    x.pboot.bar = sum(x.pboot) / n}
\NormalTok{    p.pboot.mle = r / (r + x.pboot.bar)}
\NormalTok{    p.pboot.mle}
\NormalTok{  \}}
\NormalTok{)}

\NormalTok{hist(p.pboot)}
\end{Highlighting}
\end{Shaded}

A \(95\%\) confidence interval is generated in the code block below.

\begin{Shaded}
\begin{Highlighting}[]
\NormalTok{quantile(p.pboot, c(.025, .975))}
\end{Highlighting}
\end{Shaded}

So, our approximate 95\% confidence interval for \(p\) using Parametric
Bootstrapping is about \([0.3738, 0.4336]\), depending on the
simulation.

\subsection{Nonparametric Bootstrap}\label{nonparametric-bootstrap}

for the nonparametric boostrap, we just re-sample new datasets from our
sample with replacement, and take the MLE sfrom each of these simulated
datasets. A histogram of the MLE values of these datasets can be seen in
the code block below.

\begin{Shaded}
\begin{Highlighting}[]
\NormalTok{p.npboot \textless{}{-} replicate(10000, \{}
  
\NormalTok{    \# Generate new sample (with replacement) from our above sample}
\NormalTok{    x.npboot \textless{}{-} sample(x, size=n, replace=TRUE)}
    
\NormalTok{    \# Calculate new MLE}
\NormalTok{    x.npboot.bar = sum(x.npboot) / n}
\NormalTok{    p.npboot.mle = r / (r + x.npboot.bar)}
\NormalTok{    p.npboot.mle}
\NormalTok{  \}}
\NormalTok{)}

\NormalTok{hist(p.npboot)}
\end{Highlighting}
\end{Shaded}

A \(95\%\) confidence interval is generated in the code block below.

\begin{Shaded}
\begin{Highlighting}[]
\NormalTok{quantile(p.npboot, c(.025, .975))}
\end{Highlighting}
\end{Shaded}

So, our approximate 95\% confidence interval for \(p\) using
Nonparametric Bootstrapping is about \([0.3714, 0.4360]\), depending on
the simulation.

\subsection{Comparing Confidence
Intervals}\label{comparing-confidence-intervals}

Combining last week and this week, I have now computed four different
95\% confidence intervals for \(p\).

They are listed below:

\begin{enumerate}
\def\labelenumi{\arabic{enumi}.}
\tightlist
\item
  Using CLT with \(\hat{p}_{MLE}\) plug in for \(p\) in \(I(p)\):
  \([0.3711, 0.4331]\)
\item
  Using CLT with direct computation for \(p\) in \(I(p)\):
  \([0.3727, 0.4346]\)
\item
  Parametric Bootstrap: \([0.3738, 0.4336]\) (with volatility depending
  on simulation)
\item
  Nonparametric Bootstrap: \([0.3714, 0.4360]\) (with volatility
  depending on simulation)
\end{enumerate}

Computing for the \textbf{range} of the 95\% confidence interval, we
have:

Using CLT with \(\hat{p}_{MLE}\) plug in for \(p\) in \(I(p)\):
\(0.0620\)

Using CLT with direct computation for \(p\) in \(I(p)\): \(0.0619\)

Parametric Bootstrap: \(0.0598\)

Nonparametric Bootstrap: \(0.0646\)

So, using the first simulations I ran, the parametric bootstrap's 95\%
confidence interval had the smallest range, while the nonparametric
bootstrap had the largest range. I am not sure if these mean anything,
as this may have been strictly due to a weird simulation.

\section{Metropolis Hastings}\label{metropolis-hastings}

In Chapter 1, I calculated a Bayesian posterior mean for my data using a
Beta\((1,1)\) (Uniform) prior. I solved that:

\(f(p | x) \propto p^{\alpha + nr - 1} (1-p)^{\beta + \sum x_i - 1}\).

For my uniform prior, this means that:

\(f(p | x) \propto p^{1 + nr - 1} (1-p)^{1 + \sum x_i - 1}\)

\(f(p | x) \propto p^{nr} (1-p)^{\sum x_i}\).

Now, instead of recognizing the beta distribution, let's instead run the
Metropolis Hastings Algorithm.

The algorithm will work like this. Imagine we start with a list of
length one that contains one random sample from my prior distribution
(the uniform in my case). I will then generate another sample from my
prior distribution and compare it to the last element in my list (which
in this case would just be the first generated sample).

With a certain likelihood (that I will explain in a second), I will add
this new generated sample to my list. With probability
\([1 - \text{that_likelihood}]\), I will forget about this new sample
and just add that last value in my list, to my list again. I will then
generate more and more samples and continue this process over and over.

To write out the probability from earlier, imagine \(f\) is the prior
distribution, \(g\) is the distribution we are interested in, and we are
comparing new sample \(x_j\) and previous sample \(x_{j-1}\). The
likelihood of adding generated sample \(x_j\) to the list is
\(\frac{g(x_j)}{g(x_{j-1})} \cdot \frac{f(x_{j-1})}{f(x_j)}\). The
likelihood of adding \(x_{j-1}\) again is
\(1 - \Big( \frac{g(x_j)}{g(x_{j-1})} \cdot \frac{f(x_{j-1})}{f(x_j)} \Big)\)

However, as I am using a \(\text{Uniform}(0,1)\) prior, \(f(x) = 1\) for
every \(x\). This means that the we keep a generated sample with
likelihood \(\frac{g(x_j)}{g(x_{j-1})}\), and add the older samples to
our list again with likelihood
\(\Big(1 - \frac{g(x_j)}{g(x_{j-1})} \Big)\).

After generating enough samples, and then removing the first chunk of
values from my list (called a burn-in period), the remainder of my list
should resemble the \(g\), the posterior distribution, which is the
Negative Binomial in my case.

To avoid numerical underflow, we will compute \(e^{log(f(p | x))}\).

\(f(p | x) \propto e^{ nr \log (p) + {\sum x_i} \log (1-p) + C}\).

I will generate \(15000\) samples and have a burn-in period of \(5000\).

\begin{Shaded}
\begin{Highlighting}[]

\NormalTok{\# This function calculates the posterior probability (without accounting for the constant term).}
\NormalTok{\# To avoid numerical underflow, we will first calculate the log, and then do e \^{} (log(prob))}

\NormalTok{post \textless{}{-} function(p) \{}
\NormalTok{  \# calculate log posterior (for numerical stability)}
\NormalTok{  log.p.post \textless{}{-} log(p) * (n * r) +  sum(x) * log(1{-}p)}
\NormalTok{  exp(log.p.post)}
\NormalTok{\}}

\NormalTok{\# Start with an arbitrary value of p.}
\NormalTok{p.post \textless{}{-} c(0.5)}

\NormalTok{\# Run Metropolis{-}Hastings for 15000 steps}
\NormalTok{for(i in 1:15000) \{}

\NormalTok{  p.curr \textless{}{-} p.post[length(p.post)]}
  
\NormalTok{  \# Propose a new value of p. Let\textquotesingle{}s choose from uniform(0,1). Let\textquotesingle{}s call this q.}
\NormalTok{  p.prop \textless{}{-} runif(1)}
  
\NormalTok{  \# calculate probability of choosing new\_p.}
\NormalTok{  \# This is (post(new) * q(old)) / (post(old) ( q(new)))}
\NormalTok{  prob \textless{}{-} min(1,}
\NormalTok{    (post(p.prop) * dunif(p.curr, 0, 1) /}
\NormalTok{      (post(p.curr) * dunif(p.prop, 0, 1)))}
\NormalTok{  )}

\NormalTok{  \# add new p value to list}
\NormalTok{  p.new \textless{}{-} sample(c(p.prop, p.curr),}
\NormalTok{                      size=1, prob=c(prob, 1{-}prob))}
\NormalTok{  p.post \textless{}{-} c(p.post, p.new)}
\NormalTok{\}}

\NormalTok{\# Keep the latter 10,000 samples}
\NormalTok{late.p.post \textless{}{-} p.post[5001:15000]}

\NormalTok{\# We can estimate the posterior mean by taking the mean of these samples.}
\NormalTok{mean(late.p.post)}
\end{Highlighting}
\end{Shaded}

Here, from running the code above, my Metropolis-Hastings posterior mean
is around \(0.401\), depending on the simulation.

In Chapter 1, instead of using a simulation, I recognized that my
posterior followed a \(\text{Beta}(1 + nr, 1+ \sum_{i=1}^n x_i)\)
distribution. In that situation, my posterior mean was \(0.401\).

Now, let's look at the distribution of my Metropolis-Hastings
\(\hat{p}\) in the code block below.

\begin{Shaded}
\begin{Highlighting}[]
\NormalTok{hist(late.p.post)}
\end{Highlighting}
\end{Shaded}

Also, let's generate a credible interval from my Metropolis-Hastings
data from the code block below.

\begin{Shaded}
\begin{Highlighting}[]
\NormalTok{quantile(late.p.post, c(.025, .975))}
\end{Highlighting}
\end{Shaded}

Here, my Metropolis-Hastings approximate 95\% credible interval is
around \([0.3701, 0.4283]\), depending on the simulation.

When I did this from the Beta function in Chapter 4 using the
\texttt{qbeta} function, my 95\% credible interval was
\([0.3722, 0.4328]\).

So let's further compare about these 95\% credible intervals:

\textbf{Range}:

Direct computation of the posterior distribution: \(0.0606\)

Metropolis-Hastings: \(0.0582\)

So, my approximate M-H 95\% credible interval had a slightly smaller
spread, though again, this may have just been because of a weird
simulation.

\bookmarksetup{startatroot}

\chapter{Hypothesis Testing}\label{hypothesis-testing}

\section{Background}\label{background-6}

We have previously covered solving for \(\hat{p}_{MLE}\) and
\(\hat{p}_{PM}\), evaluating them, and setting up confidence and
credible intervals. Now, we will talk about a new topic: hypothesis
testing.

This chapter is the first of three about hypothesis testing. The idea is
that we have a default hypothesis about a parameter \(\theta\). Let's
say, for example, it is that \(\theta_0 = 10\). Maybe this represents
the number of symptoms per day of medical patients.

Now, let's say we are trying out a new drug and want to see if it is
effective. So, maybe our alternative hypothesis is that
\(\theta_A = 5\), thus lowering the number of symptoms per day.
Hypothesis testing is seeing whether, based on the data, we can reject
the null hypothesis and accept the alternative.

Unlike when we solved for \(\hat{p}_{MLE}\) and \(\hat{p}_{PM}\), the
result of a hypothesis test is strictly binary. It is only where we
accept/reject the null hypothesis.

Frequentist hypothesis testing is all about the Neyman-Person lemma.
This says that if we have a distribution \(f_\theta(x)\), and
\(H_0: \theta = \theta_0\) and \(H_A: \theta = \theta_A\), then we want
to reject \(H_0\) when \(\frac{f_{\theta_A}(x)}{f_{\theta_0}(x)}\) is
maximized. This is in order to maximize the power of our hypothesis
test.

To understand why this works, let's look at hypothesis testing as a
constrained optimization problem. Theoretically we would want our test
to never make a mistake. This means we would want it to only accept the
null when \(H_0\) is true and only reject when \(H_0\) is false. In
reality, there is a strict trade-off between committing a Type I error
(rejecting \(H_0\) despite \(H_0\) being true) and maximizing power
(correctly rejecting \(H_0\)). We handle this by setting a hard cap on
our Type I error (which is referred to as \(\alpha\)), typically at
\(\alpha = 0.05\). Type I error can be very harmful, for example leading
us to deem a drug as effective even if it isn't. This is why we cap it
at \(0.05\).

Once this budget is locked in, the Neyman-Pearson lemma tells us exactly
we can maximize our power (correctly rejecting \(H_0\)) while staying
within our \(\alpha\) allocation. You can think of the likelihood ratio
\(\frac{f_{\theta_A}(x)}{f_{\theta_0}(x)}\) as a benefit-to-cost
efficiency metric for every possible data outcome. The numerator
represents the ``benefit'' (gaining power if \(H_A\) is true), while the
denominator represents the ``cost'' (eating away at our fixed \(\alpha\)
budget if \(H_0\) is true). By sorting our outcomes and choosing to
reject \(H_0\) only when this ratio is at its highest, the lemma
mathematically guarantees that we have constructed the most powerful
test possible for our chosen significance level.

In class, we showed that this principle extends into one-sided
hypothesis testing. Imagine instead of saying, \(H_A: \theta = 5\), we
said \(H_A: \theta < 10\). To solve for this, we want to find a ratio
that maximizes the power for every possible \(\theta\) below 10. What we
showed in class is that we can just choose any number below 10 (let's
say 2 for this example), and the values the maximize
\(\frac{f_{\theta_A = 2}(x)}{f_{\theta_0}(x)}\) are the same values that
maximize \(\frac{f_{\theta_A = 7}(x)}{f_{\theta_0}(x)}\) or
\(\frac{f_{\theta_A = 9}(x)}{f_{\theta_0}(x)}\), or any value below 10.
So, we can just choose one one these values below 10 to set up the ratio
for \(H_A: \theta < 10\).

Because this single ratio maximizes the power across the entire range of
possible one-sided alternative values simultaneously, the resulting test
is called a Uniformly Most Powerful (UMP) test.

\section{Uniformly Most Powerful
Test}\label{uniformly-most-powerful-test}

Let's apply this to my dataset. First, let's bring my hospital dataset
into R.

\begin{Shaded}
\begin{Highlighting}[]
\NormalTok{dat \textless{}{-} read.csv("https://raw.githubusercontent.com/Ryder4real/STATS200Q\_Book/main/transmission\_pairs\_covid\_v2.csv")}
\NormalTok{x \textless{}{-} as.vector(table(dat$infector\_id)) {-} 1}
\end{Highlighting}
\end{Shaded}

In my dataset, remember that the data that has a negative binomial
distribution is the number of people each person infected beyond the
first person (This is because all infectors in the dataset infected at
least one person, so we subtract one to have a healthy amount of zero
values). Google says that for the Flu, the average person with the
disease spreads it to 1-2 people. As our negative binomial distribution
only counts the number of additional people each person infects after
the first, then setting \(E[X] = 0.5\) feels fair. In other words, for
the flu, it reasonable to stay the average person in a hospital with a
disease infects 0.5 extra people. That would means that
\(\mathbb{E}[X] = \frac{r(1-p)}{p} = 0.5\).

Isolating for \(p\), we get:

\[
\begin{aligned}
\mathbb{E}[X] &= \frac{r(1-p)}{p} \\
\frac{p}{1-p} &= \frac{r}{\mathbb{E}[X]} \\
\frac{1-p}{p} &= \frac{\mathbb{E}[X]}{r} \\
\frac{1}{p} - 1 &= \frac{\mathbb{E}[X]}{r} \\
\frac{1}{p} &= \frac{\mathbb{E}[X] + r}{r} \\
p &= \frac{r}{\mathbb{E}[X] + r}
\end{aligned}
\]

So, we set \(\mathbb{E}[x]\) to be 0.5. Let's say for the sake of this
example that we set \(r\) to be \(0.75\). In this scenario, we get
\(p = \frac{0.75}{0.5+0.75} = 0.6\).

So, our null hypothesis is \(H_0\): \(p = 0.6\). Now, we are curious if
maybe COVID is infectious than the flu (not considering other
considerations like how my dataset is limited to people only already at
the hospital). This would mean that \(E[X]\) is greater than 0.5, and as
\(\mathbb{E}[X] = \frac{r(1-p)}{p}\), a smaller \(p\) leads to a larger
\(E[X]\).

This means that:

\(H_0\): \(p = 0.6\)

\(H_A\): \(p < 0.6\)

As mentioned in the Background section, to figure out the best X values
to reject when \(p<0.6\), all we need to do is choose one \(p<0.6\), and
the same will hold true for rest. Let's proceed with \(p=0.3\).

Again, let's proceed with \(r=0.75\).

NegBin\((r, p)\) follows the distribution
\(\binom{x + r - 1}{x} p^r (1 - p)^x\):

So, our \(H_0\) gives us NegBin\((r, p=0.6)\), which follows the
distribution:
\(f_0(x) = \prod _{i=1}^n \binom{x_i + r - 1}{x_i} 0.6^r (1 - 0.6)^{x_i}\).

Our \(H_A\) gives us NegBin\((r, p=0.3)\), which follows the
distribution:
\(f_A(x) = \prod _{i=1}^n \binom{x_i + r - 1}{x_i} 0.3^r (1 - 0.3)^{x_i}\).

We want to to maximize
\(\frac{f_A(x)}{f_0(x)} = \frac{\prod _{i=1}^n \binom{x_i + r - 1}{x_i} 0.3^r (1 - 0.3)^{x_i}}{\prod _{i=1}^n \binom{x_i + r - 1}{x_i} 0.6^r (1 - 0.6)^{x_i}}\),
as these are the \(X\) values that we maximize our
\(\frac{\text{Power}}{\alpha}\). In other words, when this value ratio
is maximized, we have the best ratio for decreasing the odds of a Type
II Error while still protecting against Type I Errors.

So, let's maximize:

\[
\begin{aligned}
\frac{f_A(x)}{f_0(x)} &= \frac{\prod _{i=1}^n \binom{x_i + r - 1}{x} 0.3^r (1 - 0.3)^{x_i}}{ \prod_{i=1}^n \binom{x_i + r - 1}{x_i} 0.6^r (1 - 0.6)^{x_i}} \\
\frac{f_A(x)}{f_0(x)} &= \frac{ \Big( \prod _{i=1}^n \binom{x_i + r - 1}{x_i} \Big) 0.3^{nr} (1 - 0.3)^{\sum x_i}}{ \Big( \prod _{i=1}^n \binom{x_i + r - 1}{x_i} \Big) 0.6^{nr} (1 - 0.6)^{\sum x_i}} \\
&= \frac{0.3^{nr} (1 - 0.3)^{\sum_{i=1}^n x_i}}{ 0.6^{nr} (1 - 0.6)^{\sum_{i=1}^n x_i}} \\
&\propto \frac{0.7^{\sum_{i=1}^n x_i}}{0.4^{\sum_{i=1}^n x_i}} \\
&= \Big(\frac{0.7}{0.4}\Big)^ {\sum_{i=1}^n x_i} \\
\end{aligned}
\]

To maximize, \(\Big(\frac{0.7}{0.4}\Big)^{\sum_{i=1}^n x_i}\), we want
to reject when \(\sum_{i=1}^n x_i\) is as large as possible.

So, we this means we set up: \(P(\sum_{i=1}^n x_i \geq c) = 0.05\).

The sum of many NegBin\((p, r)\) random variables is
NegBin\((n \cdot r, p)\), My COVID hospital dataset had 807 patients, so
\(n = 807\).

\begin{Shaded}
\begin{Highlighting}[]
\NormalTok{n = length(x)}
\NormalTok{n}
\end{Highlighting}
\end{Shaded}

So, we would use the \texttt{qnbinom} R function to set up a sum from
\(c\) to \(\infty\) over the \(H_0\) PMF that covers the upper \(5\%\)
(at most) of the distribution.

\begin{Shaded}
\begin{Highlighting}[]
\NormalTok{r = 0.75}
\NormalTok{alpha = 0.05}
\NormalTok{p = 0.6}
\NormalTok{reject\_val = qnbinom(1 {-} alpha, r * n, p)}
\NormalTok{reject\_val}
\end{Highlighting}
\end{Shaded}

We then check to see how much of the mass is in this value or above (I
write \texttt{reject\_val\ -\ 1} because we want the mass above
whichever value is inputted into the \texttt{pnbinom} function):

\begin{Shaded}
\begin{Highlighting}[]
\NormalTok{1 {-} pnbinom(reject\_val {-} 1, n * r, p)}
\end{Highlighting}
\end{Shaded}

We are above 0.05. So, we need to move our rejection bound up one value,
to take up less of the mass.

\begin{Shaded}
\begin{Highlighting}[]
\NormalTok{new\_reject\_val = reject\_val + 1}
\NormalTok{1 {-} pnbinom(new\_reject\_val {-} 1, n * r, p)}
\end{Highlighting}
\end{Shaded}

\(\alpha\) is now below \(0.05\)!

\begin{Shaded}
\begin{Highlighting}[]
\NormalTok{new\_reject\_val}
\end{Highlighting}
\end{Shaded}

So, \(c = 448\)

So, we reject if \(H_0\) \(\sum x_i \ge 448\).

Now, let's see how often we would reject in my dataset.

\begin{Shaded}
\begin{Highlighting}[]
\NormalTok{observed\_sum \textless{}{-} sum(x)}
\NormalTok{observed\_sum}
\end{Highlighting}
\end{Shaded}

As \(600 > c\) (which was \(448\)), for my COVID hospital dataset, we
reject \(H_0\). This implies that we accept the alternative hypothesis
that COVID is more contagious than the flu (though this dataset is not
perfect as it only consists of people already in the hospital who
infected at least \(1\) person, so I do not want to extrapolate too
much).

\bookmarksetup{startatroot}

\chapter{Generalized Likelihood Ratio
Test}\label{generalized-likelihood-ratio-test}

\section{Background}\label{background-7}

Last chapter, we talked about hypothesis testing and the Uniformly Most
Powerful Test. Tests like these are possible when our \(H_A\) is a
greater than or less than, like \(H_A: \theta < 1\). However, what if we
had \(H_A: \theta \neq 1\).

What maximizes power for \(\theta = 0.5\) is different than what
maximizes power for \(\theta = 1.5\). So, we cannot have Uniformly Most
Powerful Test. Instead, we set up the Generalized Likelihood Ratio Test.
The idea is still the same, where we still want to find the \(x\) values
that maximize
\(\frac{f_{\text{Alternative Hypothesis}}(x)}{f_{\text{Null Hypothesis}}(x)}\).
The problem is that if \(H_A: \theta \neq 1\), different values maximize
\(f_\text{Alternative Hypothesis}(x)\) depending on if \(\theta\) is
above or below \(1\).

So, our solution is to find a \(\theta\) that maximizes the numerator.
This means that we can re-write the ratio we want to maximize as:
\(\frac{\max_\theta \mathcal{L}(\theta)}{\mathcal{L}(\theta_0)}\).

The \(\theta\) that maximizes the likelihood is the MLE. This means that
we want to maximize the ratio
\(\frac{f_{\hat{\theta}_{MLE}}(x)}{f_{\theta_0}(x)}\).

\section{Generalized Likelihood Ratio
Test}\label{generalized-likelihood-ratio-test-1}

So, let's apply this to my dataset. First, let's bring my hospital
dataset into R.

\begin{Shaded}
\begin{Highlighting}[]
\NormalTok{dat \textless{}{-} read.csv("https://raw.githubusercontent.com/Ryder4real/STATS200Q\_Book/main/transmission\_pairs\_covid\_v2.csv")}
\NormalTok{x \textless{}{-} as.vector(table(dat$infector\_id)) {-} 1}
\end{Highlighting}
\end{Shaded}

I mention this in Chapter 6, but as a reminder:

In my hospital dataset, the data that has a negative binomial
distribution is the number of people each person infected beyond the
first person (This is because all infectors in the dataset infected at
least one person, so we subtract one to have a healthy amount of zero
values). Google says that for the Flu, the average person with the
disease spreads it to 1-2 people. As our negative binomial distribution
only counts the number of additional people each person infects after
the first, then setting \(E[X] = 0.5\) feels fair. In other words, for
the flu, it reasonable to stay the average person in a hospital with a
disease infects 0.5 extra people. That would means that
\(\mathbb{E}[X] = \frac{r(1-p)}{p} = 0.5\).

I isolated for \(p\) in Chapter 6, and got that \(p = 0.6\):

Let's say our null hypothesis is \(H_0\): \(p = 0.6\). Now, we are
curious if maybe COVID has a different degree of infectiousness than the
flu (not considering other considerations like how my dataset is limited
to people only already at the hospital). This would mean that
\(E[X] \neq 0.5\), and \(p \neq 0.6\).

So, this gives us:

\(H_0\): \(p = 0.6\)

\(H_A\): \(p \neq 0.6\)

As we mentioned above, for this type of hypothesis test, we try to
maximize the ratio:

\(\frac{\max_{p} \mathcal{L}(p)}{\mathcal{L}(p_0)}\)

We can rewrite this as:

\[
\begin{aligned}
\frac{\max_{p} \mathcal{L}(p)}{\mathcal{L}(p_0)} &= \frac{\prod _{i=1}^n \binom{x_i + r - 1}{x} p^r (1 - p)^{x_i}}{ \prod_{i=1}^n \binom{x_i + r - 1}{x_i} (p_0)^r (1 - p_0)^{x_i}} \\
\end{aligned}
\]

We know that for the numerator, the \(p\) that gives the greatest
likelihood is \(\hat{p}_{MLE}\). In Chapter 1, I calculate the
\(\hat{p}_{MLE}\) for multiple samples of the Negative Binomial
distribution is \(\frac{r}{r + \bar{x}}\).

\[
\begin{aligned}
\frac{\max_{p} \mathcal{L}(p)}{\mathcal{L}(p_0)} &= \frac{\prod _{i=1}^n \binom{x_i + r - 1}{x} (\frac{r}{r + \bar{x}})^r (1 - \frac{r}{r + \bar{x}})^{x_i}}{ \prod_{i=1}^n \binom{x_i + r - 1}{x_i} (p_0)^r (1 - p_0)^{x_i}} \\
\end{aligned}
\]

We can simplify this to:

\[
\begin{aligned}
\frac{\max_{p} \mathcal{L}(p)}{\mathcal{L}(p_0)} &= \frac{\prod _{i=1}^n \binom{x_i + r - 1}{x} (\frac{r}{r + \bar{x}})^r (1 - \frac{r}{r + \bar{x}})^{x_i}}{ \prod_{i=1}^n \binom{x_i + r - 1}{x_i} (p_0)^r (1 - p_0)^{x_i}} \\
\frac{\max_{p} \mathcal{L}(p)}{\mathcal{L}(p_0)} &= \frac{(\frac{r}{r + \bar{x}})^{nr} (1 - \frac{r}{r + \bar{x}})^{\sum_{i=1}^n x_i}}{(p_0)^{nr} (1 - p_0)^{\sum_{i=1}^n x_i}} \\
\end{aligned}
\]

We want to reject when
\(\frac{\max_{p} \mathcal{L}(p)}{\mathcal{L}(p_0)} \geq c\), where:
\[P(\texttt{test\_statistic} \geq c) = \alpha = 0.05\]

However, this does not simplify to a distribution I recognize, so I
cannot solve for \(c\).

Still, we want to maximize this function. As log() is monotonic, we can
take the logarithm of the fraction.

\[
\begin{aligned}
\log \Big(\frac{\max_{p} \mathcal{L}(p)}{\mathcal{L}(p_0)} \Big) &= \log \Bigg( \frac{(\frac{r}{r + \bar{x}})^{nr} (1 - \frac{r}{r + \bar{x}})^{\sum_{i=1}^n x_i}}{(p_0)^{nr} (1 - p_0)^{\sum_{i=1}^n x_i}} \Bigg)\\
\log \Big(\frac{\max_{p} \mathcal{L}(p)}{\mathcal{L}(p_0)} \Big) &= nr \log(\frac{r}{r + \bar{x}}) + \sum_{i=1}^n x_i \log(1 - \frac{r}{r + \bar{x}})- \Big( nr \log(p_0)  + \sum_{i=1}^n x_i \log(1 - p_0) \Big)\\
\log \Big(\frac{\max_{p} \mathcal{L}(p)}{\mathcal{L}(p_0)} \Big) &= nr \log \Big(\frac{\frac{r}{r + \bar{x}}}{p_0}\Big) + \sum_{i=1}^n x_i \log \Big( \frac{1-\frac{r}{r + \bar{x}}}{1-p_0} \Big) \\
\log \Big(\frac{\max_{p} \mathcal{L}(p)}{\mathcal{L}(p_0)} \Big) &= nr \log \Big(\frac{\frac{r}{r + \bar{x}}}{p_0}\Big) + \sum_{i=1}^n x_i \log \Big( \frac{1-\frac{r}{r + \bar{x}}}{1-p_0} \Big)
\end{aligned}
\]

Once again, this also does not simplify to a distribution I recognize.
So, we are left to either simulate or take the asymptotic distribution.

\subsection{Simulation}\label{simulation}

Let's simulate many times under the null hypothesis, and calculate
\(\frac{\max_{p} \mathcal{L}(p)}{\mathcal{L}(p_0)}\) each time. For
brevity, I will refer to
\(\frac{\max_{p} \mathcal{L}(p)}{\mathcal{L}(p_0)}\) as \(\Lambda\).
Then, as we want to keep \(\alpha \leq 0.05\), we would reject if
\(\Lambda\) for our data is greater than the \(95\)th percentile of the
simulated \(\Lambda\)s.

Remember:
\(\Lambda = \frac{\max_{p} \mathcal{L}(p)}{\mathcal{L}(p_0)} = \frac{(\frac{r}{r + \bar{x}})^{nr} (1 - \frac{r}{r + \bar{x}})^{\sum_{i=1}^n x_i}}{(p_0)^{nr} (1 - p_0)^{\sum_{i=1}^n x_i}}\).

However, to prevent underflow, I will take the logarithm.

\(\log \Big(\frac{\max_{p} \mathcal{L}(p)}{\mathcal{L}(p_0)} \Big) = nr \log \Big(\frac{\frac{r}{r + \bar{x}}}{p_0}\Big) + \sum_{i=1}^n x_i \log \Big( \frac{1-\frac{r}{r + \bar{x}}}{1-p_0} \Big)\)

\begin{Shaded}
\begin{Highlighting}[]
\NormalTok{n = length(x)}
\NormalTok{r = 0.75}
\NormalTok{p\_0 = 0.6}
\NormalTok{n\_sims = 10000}

\NormalTok{sims \textless{}{-} replicate(n\_sims, \{}
\NormalTok{    data \textless{}{-} rnbinom(n = n, size = r, prob = p\_0)}
    
\NormalTok{    data\_bar = mean(data)}
\NormalTok{    data\_sum = sum(data)}
    
\NormalTok{    p\_mle = r / (r + data\_bar)}
    
\NormalTok{    first\_term = n * r * log(p\_mle / p\_0)}
\NormalTok{    second\_term = data\_sum * log((1 {-} p\_mle) / (1 {-} p\_0))}
    
\NormalTok{    log\_lambda = first\_term + second\_term}
\NormalTok{    log\_lambda}
\NormalTok{\}}
\NormalTok{)}
\NormalTok{quantile(sims, 0.95)}
\end{Highlighting}
\end{Shaded}

So, our test statistic is
\(\log \Big(\frac{\max_{p} \mathcal{L}(p)}{\mathcal{L}(p_0)} \Big)\) and
our \(c\) is about \(1.933\), depending on the simulation. So, if
\(\Lambda\) for my data is above 1.933, we reject \(H_0\).

So, let's calculate \(\log (\Lambda)\), while setting \(r=0.75\).

\begin{Shaded}
\begin{Highlighting}[]
\NormalTok{n = length(x)}
\NormalTok{r = 0.75}
\NormalTok{p\_0 = 0.6}

\NormalTok{x\_bar = mean(x)}
\NormalTok{x\_sum = sum(x)}
    
\NormalTok{p\_mle = r / (r + x\_bar)}
    
\NormalTok{first\_term = n * r * log(p\_mle / p\_0)}
\NormalTok{second\_term = x\_sum * log((1 {-} p\_mle) / (1 {-} p\_0))}
    
\NormalTok{log\_lambda = first\_term + second\_term}
\NormalTok{log\_lambda}
\end{Highlighting}
\end{Shaded}

As \(\log (\Lambda)\) is 23.55, which is \(>c\), we reject \(H_0\).

A power curve stricly from simulation is generated in the code block
below.

Note, for the sake of runtime, I lowered \texttt{n} to 50 in the code. I
recommend messing around and changing \texttt{n} to see how the power
curve changes.

\begin{Shaded}
\begin{Highlighting}[]
\NormalTok{\# I encourage you to change n to see how this affects the power curve}
\NormalTok{n \textless{}{-} 50}

\NormalTok{p\_0 \textless{}{-} 0.6}
\NormalTok{r = 0.75}
\NormalTok{n\_sims \textless{}{-} 5000  }

\NormalTok{\# Simulating under the null hypothesis to find the 95th quartile (our c value)}
\NormalTok{null\_sims \textless{}{-} replicate(n\_sims, \{}
\NormalTok{  data \textless{}{-} rnbinom(n = n, size = r, prob = p\_0)}
\NormalTok{  data\_bar \textless{}{-} mean(data)}
  
\NormalTok{  p\_mle \textless{}{-} r / (r + data\_bar)}
  
\NormalTok{  \# Avoid any potential log(0) errors}
\NormalTok{  if (p\_mle == 1) \{ }
\NormalTok{    p\_mle \textless{}{-} 0.9999 }
\NormalTok{  \}}
  
\NormalTok{  \# Log{-}likelihood ratio logic}
\NormalTok{  log\_lambda \textless{}{-} n * r * log(p\_mle / p\_0) + sum(data) * log((1 {-} p\_mle) / (1 {-} p\_0))}
\NormalTok{  log\_lambda}
\NormalTok{\})}
\NormalTok{raw\_crit\_val \textless{}{-} quantile(null\_sims, 0.95)}

\NormalTok{\# Set up parameters for power curve}
\NormalTok{p\_alt\_space \textless{}{-} seq(0.01, 1, by = 0.05)}
\NormalTok{powers \textless{}{-} c()  }

\NormalTok{\# Simulate alternative hypothesis / null hypothesis and keep track for each one whether we should reject}

\NormalTok{for (p\_alt in p\_alt\_space) \{}
  
\NormalTok{  rejections \textless{}{-} replicate(n\_sims, \{}

\NormalTok{    data\_alt \textless{}{-} rnbinom(n = n, size = r, prob = p\_alt)}
\NormalTok{    data\_bar\_alt \textless{}{-} mean(data\_alt)}
    
\NormalTok{    p\_mle\_alt \textless{}{-} r / (r + data\_bar\_alt)}
    
\NormalTok{    \# Avoid any potential log(0) errors}
\NormalTok{    if (p\_mle\_alt == 1) \{ }
\NormalTok{      p\_mle\_alt \textless{}{-} 0.9999 }
\NormalTok{    \}}
    
\NormalTok{    \# Calculate log\_lambda}
\NormalTok{    log\_lambda\_alt \textless{}{-} n * r * log(p\_mle\_alt / p\_0) + sum(data\_alt) * log((1 {-} p\_mle\_alt) / (1 {-} p\_0))}
    
\NormalTok{    \# Check if we should reject?}
\NormalTok{    log\_lambda\_alt \textgreater{} raw\_crit\_val}
\NormalTok{  \})}
  
\NormalTok{  \# Append the mean rejection rate (power) for this specific p\_alt}
\NormalTok{  powers \textless{}{-} c(powers, mean(rejections))}
\NormalTok{\}}

\NormalTok{\# Plot results}
\NormalTok{plot(p\_alt\_space, powers, type = "l", lwd = 3, col = "\#2c3e50",}
\NormalTok{     xlab = expression(p), ylab = "Power", }
\NormalTok{     main = "Empirical GLRT Power Curve", ylim = c(0, 1))}

\NormalTok{abline(v = p\_0, col = "\#e74c3c", lty = 2)  }
\NormalTok{abline(h = 0.05, col = "gray50", lty = 3) }
\end{Highlighting}
\end{Shaded}

I recommend going back in the code and seeing how changing \texttt{n}
can effect your power curve.

\subsection{Asymptotic}\label{asymptotic}

We learned in class that under Wilks' Theorem, for large \(n\):

\(2 \log(\frac{\max_{p} \mathcal{L}(p)}{\mathcal{L}(p_0)}) \sim \chi^2_1\).

As we want to reject when Remember, we want to reject when
\(2 \log(\frac{\max_{p} \mathcal{L}(p)}{\mathcal{L}(p_0)})\) is large.

So, we set up our rejection bound as:

\(2 \log(\frac{\max_{p} \mathcal{L}(p)}{\mathcal{L}(p_0)}) \geq c\),
where \(c\) the \(95\%\) percentile of the chi-squared distribution.

So, let's figure out our \(c\):

We can run this in code:

\begin{Shaded}
\begin{Highlighting}[]
\NormalTok{alpha = 0.05}
\NormalTok{p = 0.6}
\NormalTok{c = qchisq(1{-}alpha, 1)}
\NormalTok{c}
\end{Highlighting}
\end{Shaded}

So, form the code block above, \(c = 3.841459\)

Now, let's check our data. Let's assume \(r = 0.75\). Remember that
Wilks' Theorem states that for large \(n\):
\(2 \log(\frac{\max_{p} \mathcal{L}(p)}{\mathcal{L}(p_0)}) \sim \chi^2_1\).

For our Negative Binomial:
\(\log \Big(\frac{\max_{p} \mathcal{L}(p)}{\mathcal{L}(p_0)} \Big) = nr \log \Big(\frac{\frac{r}{r + \bar{x}}}{p_0}\Big) + \sum_{i=1}^n x_i \log \Big( \frac{1-\frac{r}{r + \bar{x}}}{1-p_0} \Big)\)

Let's proceed with \(r=0.75\).

\begin{Shaded}
\begin{Highlighting}[]
\NormalTok{n = length(x)}
\NormalTok{r = 0.75}
\NormalTok{x\_sum = sum(x)}
\NormalTok{x\_bar = mean(x)}
\NormalTok{p\_0 = 0.6}
\NormalTok{p\_mle = r / (r + x\_bar)}
\NormalTok{log\_part = n * r * log(p\_mle / p\_0) + x\_sum * log((1 {-} p\_mle) / (1 {-} p\_0))}
\NormalTok{test\_stat = 2 * log\_part}
\NormalTok{test\_stat}
\end{Highlighting}
\end{Shaded}

So, we can see from the code block that our dataset's test statistic was
47.09489. If I just plug in the PDFs, I get the same value:

\begin{Shaded}
\begin{Highlighting}[]
\NormalTok{log\_L\_mle  \textless{}{-} dnbinom(x\_sum, size = n * r, prob = p\_mle, log = TRUE)}
\NormalTok{log\_L\_null \textless{}{-} dnbinom(x\_sum, size = n * r, prob = p\_0, log = TRUE)}

\NormalTok{test\_stat\_2 = 2 * (log\_L\_mle {-} log\_L\_null)}
\NormalTok{test\_stat\_2}
\end{Highlighting}
\end{Shaded}

As \(47.1 > c\) (which was \(3.841459\)), for my COVID hospital dataset,
we reject \(H_0\).

This implies that we reject the hypothesis that COVID has the same
contagious level as the flu (though this dataset is not perfect as it
only consists of people already in the hospital who infected at least
\(1\) person, so I do not want to extrapolate too much). Another way to
say this is that we accept the hypothesis that COVID does not have the
same contagious level as the flu.

A power curve using Wilk's Theorem would be:

\begin{Shaded}
\begin{Highlighting}[]
\NormalTok{\# Feel free to change this parameter (n) to see how it affects the curve}
\NormalTok{n = length(x)}

\NormalTok{p\_alt \textless{}{-} seq(0, 1, by = 0.001)}
\NormalTok{alpha \textless{}{-} 0.05}
\NormalTok{crit\_val \textless{}{-} qchisq(alpha, df = 1, lower.tail = FALSE)}

\NormalTok{lambda \textless{}{-} 2 * n * r * (log(p\_alt / p\_0) + ((1 {-} p\_alt) / p\_alt) * log((1 {-} p\_alt) / (1 {-} p\_0)))}

\NormalTok{lambda[p\_alt == p\_0] \textless{}{-} 0}

\NormalTok{power\_curve \textless{}{-} pchisq(crit\_val, df = 1, ncp = lambda, lower.tail = FALSE)}

\NormalTok{plot(p\_alt, power\_curve, type = "l", col = "\#2c3e50", lwd = 3,}
\NormalTok{     xlab = expression(textstyle("p")),}
\NormalTok{     ylab = "Power",}
\NormalTok{     main = "GLR Power Curve (Wilks\textquotesingle{} Theorem)",}
\NormalTok{     ylim = c(0, 1), panel.first = grid(lty = "dotted", col = "gray70"))}

\NormalTok{abline(v = p\_0, col = "\#e74c3c", lty = 2, lwd = 2)   \# Null Hypothesis baseline}

\NormalTok{legend("bottomright", }
\NormalTok{       legend = c("Power Curve", "Null Hypothesis (p = 0.6)"),}
\NormalTok{       col = c("\#2c3e50", "\#e74c3c", "\#27ae60"), }
\NormalTok{       cex=0.6,}
\NormalTok{       lty = c(1, 2, 3), lwd = c(3, 2, 2), bg = "white")}
\end{Highlighting}
\end{Shaded}

Again, I recommend going back in the code and seeing how changing
\texttt{n} can effect your power curve.

\bookmarksetup{startatroot}

\chapter{Bayesian Hypothesis Testing}\label{bayesian-hypothesis-testing}

\section{Background}\label{background-8}

In the previous two chapters, we have talked about frequentist
hypothesis testing. We now turn our heads to Bayesian hypothesis
testing. For Bayesian hypothesis testing, we can compare the posterior
likelihoods of \(H_0\) and \(H_A\). This is called the Bayes Factor. If
the Bayes Factor is above 20, then we are considered to have ``strong
evidence'' to reject the null.

Similar to our previous Bayesian metrics, we first need some priors.
There are three priors we will need to come up with. The first two are
the prior likelihoods of \(H_0\) and \(H_A\), which naturally need to
sum to 1. The third is a prior distribution for \(\theta\).

\section{Bayesian Hypothesis Test}\label{bayesian-hypothesis-test}

First, let's bring my hospital dataset into R.

\begin{Shaded}
\begin{Highlighting}[]
\NormalTok{dat \textless{}{-} read.csv("https://raw.githubusercontent.com/Ryder4real/STATS200Q\_Book/main/transmission\_pairs\_covid\_v2.csv")}
\NormalTok{x \textless{}{-} as.vector(table(dat$infector\_id)) {-} 1}
\end{Highlighting}
\end{Shaded}

I mention this in Chapter 6, but as a reminder:

In my hospital dataset, the data that has a negative binomial
distribution is the number of people each person infected beyond the
first person (This is because all infectors in the dataset infected at
least one person, so we subtract one to have a healthy amount of zero
values). Google says that for the Flu, the average person with the
disease spreads it to 1-2 people. As our negative binomial distribution
only counts the number of additional people each person infects after
the first, then setting \(E[X] = 0.5\) feels fair. In other words, for
the flu, it reasonable to stay the average person in a hospital with a
disease infects 0.5 extra people. That would means that
\(\mathbb{E}[X] = \frac{r(1-p)}{p} = 0.5\).

I isolated for \(p\) in Chapter 6, and got that \(p = 0.6\):

Let's say our null hypothesis is \(H_0\): \(p = 0.6\). Now, we are
curious if maybe COVID has a different degree of infectiousness than the
flu (not considering other considerations like how my dataset is limited
to people only already at the hospital). This would mean that
\(E[X] \neq 0.5\), and and \(p \neq 0.6\).

So, this gives us:

\(H_0\): \(p = 0.6\)

\(H_A\): \(p \neq 0.6\)

Now, for Bayesian Hypothesis Testing, we need to assign prior
probabilities to both the null hypothesis and the alternative
hypothesis. As I do not have a great intuition, let's follow the Rule of
Insufficient Reason and say both are \(0.5\).

So, we have:

\(P(H_0) = 0.5\)

\(P(H_A) = 0.5\)

Next, we need to find the prior distribution for \(p\) in \(H_A\).
Again, as I do not have a great intuition, let's once again follow the
Rule of Insufficient Reason and say \(p \sim \text{Uniform}(0,1)\), as
\(p\) could be anything from \(0\) to \(1\).

So, using Bayes' Rule, this means that:

\(P(H_A|x) = \frac{P(H_A) f(x|H_A)}{f(x)}\)

\(P(H_0|x) = \frac{P(H_0) f(x|H_0)}{f(x)}\)

NOTE: In the formula below, I will refer to \(p\) in \(H_0\) as \(p_0\).
I know we set this equal to \(0.6\), but for clarity, I think it is
easier to understand with \(p_0\).

Remember that \(f(x_1,x_2,...x_n)\) is the product of Negative Binomial
random variables.

When taking their ratio, we get the Bayes Factor, which is:

\[
\begin{aligned}
\frac{P(H_A|x_1,...,x_n)}{P(H_0|x_1,...,x_n)} &= \frac{\frac{P(H_A) f(x_1,...,x_n|H_A)}{f(x_1,...,x_n)}}{\frac{P(H_0) f(x_1,...,x_n|H_0)}{f(x_1,...,x_n)}} \\
&= \frac{P(H_A) f(x_1,...,x_n|H_A)}{P(H_0) f(x_1,...,x_n|H_0)} \\
&= \frac{P(H_A)}{P(H_0)} \frac{f(x_1,...,x_n|H_A)}{f(x_1,...,x_n|H_0)} \\
&= \frac{0.5}{0.5} \frac{\int _0^1 f(p) f(x_1,...,x_n|p)dp}{\prod_{i=1}^n \binom{x_i + r - 1}{x_i} (p_0)^r (1 - p_0)^{x_i}} \\
&= \frac{\int _0^1 f(x_1,...,x_n|p)dp}{\prod_{i=1}^n \binom{x + r - 1}{x} (p_0)^r (1 - p_0)^{x_i}} \\
&= \frac{\int _0^1 \prod_{i=1}^n \binom{x_i + r - 1}{x_i} p^r (1 - p)^{x_i} dp}{\prod_{i=1}^n \binom{x_i + r - 1}{x_i} (p_0)^r (1 - p_0)^{x_i}} \\
&= \frac{\int_0^1 \prod_{i=1}^n \binom{x_i + r - 1}{x_i} p^r (1 - p)^{x_i} \, dp}{\prod_{i=1}^n \binom{x_i + r - 1}{x_i} (p_0)^r (1 - p_0)^{x_i}} \\
&= \frac{\int_0^1 \left( \prod_{i=1}^n \binom{x_i + r - 1}{x_i} \right) p^{rn} (1 - p)^{\sum_{i=1}^n x_i} \, dp}{\left( \prod_{i=1}^n \binom{x_i + r - 1}{x_i} \right) (p_0)^{rn} (1 - p_0)^{\sum_{i=1}^n x_i}} \\
&= \frac{ \left( \prod_{i=1}^n \binom{x_i + r - 1}{x_i} \right) \int_0^1 p^{rn} (1 - p)^{\sum_{i=1}^n x_i} \, dp}{\left( \prod_{i=1}^n \binom{x_i + r - 1}{x_i} \right) (p_0)^{rn} (1 - p_0)^{\sum_{i=1}^n x_i}} \\
&= \frac{\int_0^1 p^{rn} (1 - p)^{\sum_{i=1}^n x_i} \, dp}{(p_0)^{rn} (1 - p_0)^{\sum_{i=1}^n x_i}}
\end{aligned}
\]

We can simplify the numerator by recognizing that it has the meat of the
Beta PDF over its domain.

\[
\begin{aligned}
\frac{P(H_A | x)}{P(H_0 | x)} &= \frac{\int_0^1 p^{rn} (1 - p)^{\sum_{i=1}^n x_i} \, dp}{(p_0)^{rn} (1 - p_0)^{\sum_{i=1}^n x_i}} \\
&= \frac{\frac{\Gamma(rn + 1) \Gamma(\sum_{i=1}^n x_i + 1)}{\Gamma(rn + 1 + \sum_{i=1}^n x_i + 1) } \int_0^1 \frac{\Gamma(rn + 1 + \sum_{i=1}^n x_i + 1)}{\Gamma(rn + 1) \Gamma(\sum_{i=1}^n x_i + 1)}p^{rn} (1 - p)^{\sum_{i=1}^n x_i} \, dp}{(p_0)^{rn} (1 - p_0)^{\sum_{i=1}^n x_i}} \\
&= \frac{\frac{\Gamma(rn + 1) \Gamma\left(\sum_{i=1}^n x_i + 1\right)}{\Gamma\left(rn + \sum_{i=1}^n x_i + 2\right)}}{(p_0)^{rn} (1 - p_0)^{\sum_{i=1}^n x_i}}
\end{aligned}
\]

In my hospital dataset, the data seems to fit a negative binomial
distribution in which \(r\) is not necessarily an integer, meaning we
cannot simplify this further with factorials.

Instead, let's plug our data this equation to find the Bayes Factor.
Let's say \(r = 0.75\), the same \(r\) value I used in Chapter 7.

Unfortunately, the denominator gets so small that the code thinks we are
dividing by \(0\). So, let's instead take \(e\) to the power of {[}the
\(\log\) of the numerator subtracted by the \(log\) of the
denominator{]}.

\[
\begin{aligned}
\frac{P(H_A | x)}{P(H_0 | x)}
&= \frac{\frac{\Gamma(rn + 1) \Gamma\left(\sum_{i=1}^n x_i + 1\right)}{\Gamma\left(rn + \sum_{i=1}^n x_i + 2\right)}}{(p_0)^{rn} (1 - p_0)^{\sum_{i=1}^n x_i}} \\
&= e^{\log \Bigg( \frac{\frac{\Gamma(rn + 1) \Gamma\left(\sum_{i=1}^n x_i + 1\right)}{\Gamma\left(rn + \sum_{i=1}^n x_i + 2\right)}}{(p_0)^{rn} (1 - p_0)^{\sum_{i=1}^n x_i}} \Bigg)} \\
&= e^{\log \Big( \frac{\Gamma(rn + 1) \Gamma\left(\sum_{i=1}^n x_i + 1\right)}{\Gamma\left(rn + \sum_{i=1}^n x_i + 2\right)} \Big) - \log\Big( (p_0)^{rn} (1 - p_0)^{\sum_{i=1}^n x_i} \Big)} \\
&= e^{\log \Big(\Gamma(rn + 1) \Big) + \log \Big( \Gamma\left(\sum_{i=1}^n x_i + 1\right)\Big) - \log \Big(\Gamma\left(rn + \sum_{i=1}^n x_i + 2\right) \Big) - \Big(rn\log(p_0) + {\sum_{i=1}^n x_i}\log(1 - p_0) \Big)}
\end{aligned}
\]

\begin{Shaded}
\begin{Highlighting}[]
\NormalTok{r = 0.75}
\NormalTok{sum\_x = sum(x)}
\NormalTok{n = length(x)}
\NormalTok{p\_0 = 0.6}

\NormalTok{first\_term = lgamma(r * n+1)}
\NormalTok{second\_term = lgamma(sum\_x + 1)}
\NormalTok{third\_term = lgamma(r * n + sum\_x + 2)}
\NormalTok{fourth\_term = r * n * log(p\_0) + sum\_x * log(1{-}p\_0)}

\NormalTok{exponent = first\_term + second\_term {-} third\_term {-} fourth\_term}

\NormalTok{bayes\_factor = exp(exponent)}
\NormalTok{bayes\_factor}
\end{Highlighting}
\end{Shaded}

My Bayes Factor is \(607815364\), meaning there is very strong evidence
to reject the null hypothesis per the criteria in Kass and Raftery
(1995).

In Chapter 7, using Frequentist Hypothesis Testing methods, I also
rejected the null hypothesis. So, both methods reached the same
conclusion, in that we rejected \(H_0\) (we reject the hypothesis that
COVID has the same contagious level as the flu). As my Bayesian and
frequenist hypothesis tests both come to the same conclusion, I did not
encounter Lindley's paradox.

\bookmarksetup{startatroot}

\chapter{Sufficency}\label{sufficency}

\section{Background}\label{background-9}

We have finally finished covering both frequentist and Bayesian
hypothesis testing. This means we moved onto the final topic of this
book: sufficiency.

A sufficient statistic \(T(x_1,...,x_n)\), means that \(T\) contains all
useful information with respect to the data. This means that once we
have \(T\), we do not care how we got there. As a quick example, let's
say that \(T = \sum_{i=1}^n x_i\) and it is a sufficient statistic for
\(\theta\). This means that for properly estimating \(\theta\), we do
not need to know \(x_1,...,x_n\). In fact, we only need to know \(T\).

Furthermore, for a frequentist application, the Rao-Blackwell theorem
says that for any unbiased estimator \(\hat{\theta}\), then
\(\hat{\theta}_{RB} = \mathbb{E}[\hat{\theta}(T)]\) is also unbiased and
\(\text{Var}(\hat{\theta}_{RB}) \leq \text{Var}(\hat{\theta})\). So,
this means that the best unbiased estimator only references the data
with respect to the sufficient statistic \(T\).

A quick proof is below. First, let's show that if \(\hat{\theta}\) is
unbiased, then so is \(\hat{\theta}_{RB}\).

\[
\begin{aligned}
  \mathbb{E}[\hat{\theta}_{RB}] &= \mathbb{E}[\mathbb{E}[\hat{\theta} | T]]  \\
  \mathbb{E}[\hat{\theta}_{RB}] &= \mathbb{E}[\hat{\theta}] \text{   [NOTE: From the Law of Total Exectation]}
\end{aligned}
\]

So, this means that the expected value of both estimators is the same,
so they have the same bias. Next, let's show the variance part.

\[
\begin{aligned}
\text{Var}(\hat{\theta}) &= \mathbb{E}[(\text{Var}(\hat{\theta}|T)] + \text{Var}(\mathbb{E}[\hat{\theta}|T]) \text{   [NOTE: From the Law of Total Variance]} \\
\text{Var}(\hat{\theta}) &= (\text{value} \geq 0) + \text{Var}(\hat{\theta}_{RB})
\end{aligned}
\]

So, we can see
\(\text{Var}(\hat{\theta}_{RB}) \leq \text{Var}(\hat{\theta})\).

This gives way to the sufficiency principle: ``If you have a sufficient
statistic, use it''.

The MLE and Posterior Mean are always both functions of a sufficient
statistic. In general, our estimators are usually already functions of
sufficient statistics, meaning we do not usually have to apply the
Rao-Blackwell theorem to improve our estimator (this process is called
Rao-Blackwellization).

\section{Sufficient Statistic}\label{sufficient-statistic}

In class, we learned that for data \(X_i,...X_n\) and parameter
\(\theta\), a statistic \(T(x_i,...x_n)\) is sufficient if either of the
following hold:

1:
\(f_\theta(x_i,...x_n) = g\Big( \theta, T(x_1,...,x_n)\Big) \cdot h(x_1,,...x_n)\),

\(\quad\) where:

\(\qquad g\Big( \theta, T(x_1,...,x_n)\Big)\) depends on \(\theta\) only
through \(T\)

\(\qquad h(x_1,,...x_n)\) depends only on \((x_1,,...x_n)\) and not on
\(\theta\)

\begin{enumerate}
\def\labelenumi{\arabic{enumi}.}
\setcounter{enumi}{1}
\tightlist
\item
  \(P(X_i=x_1,...,X_n=x_n | T = t)\) does not depend on \(\theta\).
\end{enumerate}

For my Negative Binomial dataset, \(p\) is the parameter we are
interested in estimating, so \(p\) will be our \(\theta\).

I will show that for my distribution, both definitions hold:

\subsection{First Proof}\label{first-proof}

The PDF of my negative Binomial Distribution:
\(X_1,...,X_n \sim \text{NegBin}(r, p)\) is:

\[
\begin{aligned}
    f_p(x_1,...,x_n) &= \prod_{i=1}^n \binom{x_i + r- 1}{x_i} p^r (1 - p)^{x_i}
\end{aligned}
\]

We can simplify this into:

\[
\begin{aligned}
    f_p(x_1,...,x_n) &= \prod_{i=1}^n \binom{x_i + r- 1}{x_i} p^r (1 - p)^{x_i} \\
     &= p^{rn} \prod_{i=1}^n \binom{x_i + r- 1}{x_i} (1 - p)^{x_i} \\
    &= p^{rn} \Bigg( \prod_{i=1}^n \binom{x_i + r- 1}{x_i} \Bigg) (1 - p)^{\sum_{i=1}^nx_i} 
\end{aligned}
\]

We can re-write this as:

\[
\begin{aligned}
 f_p(x_1,...,x_n) &= p^{rn}(1 - p)^{\sum_{i=1}^n x_i} \cdot  \prod_{i=1}^n \binom{x_i + r- 1}{x_i} 
\end{aligned}
\]

Now, as we have been doing, let's assume \(r\) is fixed and \(p\) is the
parameter we are interested in estimating.

Here, if we set \(T(X_1,...x_n) = \sum_{i=1}^n x_i\), then we can
rewrite the equation as:

\[
\begin{aligned}
 f_p(x_1,...,x_n) &= p^{rn}(1 - p)^{\sum_{i=1}^n x_i} \cdot  \prod_{i=1}^n \binom{x_i + r- 1}{x_i} \\
&= p^{rn}(1 - p)^{T} \cdot  \prod_{i=1}^n \binom{x_i + r- 1}{x_i} \\
&= g\Big( p, T(x_1,...,x_n)\Big)  \cdot  h(x_1,,...x_n)
\end{aligned}
\]

where:

\begin{enumerate}
\def\labelenumi{\arabic{enumi}.}
\item
  \(g\Big( p, T(x_1,...,x_n)\Big) = p^{rn}(1 - p)^{T}\)
\item
  \(h(x_1,,...x_n) = \prod_{i=1}^n \binom{x_i + r- 1}{x_i}\)
\end{enumerate}

As we were able to write our PDF in terms of
\(g\Big(p, T(x_1,...,x_n)\Big)  \cdot  h(x_1,,...x_n)\), then this means
\(T(x_1,...,x_n)\) is a sufficient statistic.

So, for the Negative Binomial distribution, \(\sum_{i=1}^n x_i\) is a
sufficient statistic.

\subsection{Second Proof}\label{second-proof}

Similarly, we could have also confirmed this using the other definition
we learned in class, that \(P(X_1=x_1,...,X_n=x_n | T = t)\) does not
depend on \(p\).

\[
\begin{aligned}
P(X_1=x_1,...,Xn=x_n | T=t) = \frac{P(X_1=x_1,...,X_n=x_n, T=t)}{P(T=t)} 
\end{aligned}
\]

\(T\) is the sum of i.i.d. \(\text{NegBin}(r, p)\) random variables. So,
\(T\) is a \(\text{NegBin}(nr, p)\) random variable.

So, we get:

\[
\begin{aligned}
P(X_1=x_1,...,Xn=x_n | T=t) &= \frac{P(X_1=x_1,...,X_n=x_n, T=t)}{P(T=t)} \\
P(X_1=x_1,...,Xn=x_n | T=t) &= \frac{p^{rn}(1 - p)^{t} \cdot  \prod_{i=1}^n \binom{x_i + r - 1}{x_i}}{\binom{t + nr- 1}{t} p^{rn} (1 - p)^{t}} \\
P(X_1=x_1,...,Xn=x_n | T=t) &= \frac{  \prod_{i=1}^n \binom{x_i + r - 1}{x_i}}{\binom{t + nr - 1}{t} } \\
\end{aligned}
\]

As we can see, this does not depend on \(p\), proving once again that
\(T(x_1,...x_n) = \sum_{i=1}^n x_i\) is a sufficient statistic.

\section{Seeing This in My MLE and Posterior
Mean}\label{seeing-this-in-my-mle-and-posterior-mean}

As I mentioned earlier, the MLE and Posterior Mean are both always
functions of a Sufficient Statistic. From Chapter 1, we have that:

\[
\begin{aligned}
\hat{p}_{MLE} &= \frac{r}{r + \bar{x}} \\
\hat{p}_{PM} &= \frac{1 + nr}{2 + nr + \sum_{i=1}^n x_i}
\end{aligned}
\]

As \(T(x_1,...x_n) = \sum_{i=1}^n x_i\), we can rewrite our MLE and
Posterior Mean as functions of T:

\[
\begin{aligned}
\hat{p}_{MLE} &= \frac{r}{r + \frac{T}{n}} \\
\hat{p}_{posterior} &= \frac{1 + nr}{2 + nr + T}
\end{aligned}
\]

\bookmarksetup{startatroot}

\chapter*{References}\label{references}
\addcontentsline{toc}{chapter}{References}

\markboth{References}{References}

\protect\phantomsection\label{refs}
\begin{CSLReferences}{1}{1}
\bibitem[\citeproctext]{ref-lloydsmith2005}
Lloyd-Smith, Jamie O., Sebastian J. Schreiber, P. Ekkehard Kopp, and
Wayne M. Getz. 2005. {``Superspreading and the Effect of Individual
Variation on Disease Emergence.''} \emph{Nature} 438 (7066): 355--59.
\url{https://doi.org/10.1038/nature04153}.

\end{CSLReferences}




\end{document}
